\begin{appendixs}
	
% \section{Modelos de movimiento de suelo}
% \label{anexo: Modelos de movimiento de suelo}

% \subsection{Subducción}

% \subsubsection{Modelos de movimiento del suelo para fuentes interplaca e intraplaca}
% \label{sec:modelos_inter_intra}



% Los modelos de NGA-Sub (AG, KBCG y PSBAH) fueron desarrollados utilizando la misma base de datos (\textit{NGA-Sub database}), que contiene información de siete regiones tectónicas de subducción: Alaska, América Central y México, Cascadia, Japón, Nueva Zelanda, Sudamérica y Taiwán. Aunque estos modelos comparten la base de datos, los criterios de selección y las decisiones de modelación generan diferencias en sus predicciones \shortcite{Gregor2022}. El modelo BCHU \shortcite{Abrahamson2018} también emplea esta base de datos.  

% El modelo MBR \shortcite{Montalva2017} se basa en 3774 registros de 473 eventos ocurridos entre 1985 y 2015 en Chile. Utiliza la forma funcional desarrollada por Abrahamson et al.\ (2016) para el proyecto BC Hydro, incorporada a partir de dicho modelo.

% En la tabla \ref{tab:Parametros de aplicabilidad GMPE SUB} se muestran los rangos de valores en donde son aplicables los modelos de movimiento de suelo. Del mismo modo, en la tabal \ref{tab:Variables usadas GMPE SUB} se muestran las variables de entrada utilizadas para cada modelo.

% \subsubsection{Incertidumbre epistémica y variabilidad aleatoria}
% \label{sec:incertidumbre_y_variabilidad}

% \subsubsection{Incertidumbre epistémica}
% La incertidumbre epistémica está asociada con la falta de conocimiento sobre los procesos de generación del movimiento del suelo. Cada GMM maneja esta incertidumbre de forma particular:

% \begin{itemize}
%   \item \textbf{Montalva et al.\ (2017).}  
%   Emplea la forma funcional y varios coeficientes del modelo de Abrahamson et al.\ (2016), también denominado modelo BC Hydro. La incertidumbre epistémica de dicho modelo se transfiere al de Montalva et al. Además, se desarrolla un modelo alternativo basado en datos de alta calidad (HQ) para analizar el impacto de la calidad de los datos.

%   \item \textbf{Parker et al.\ (2020).}  
%   Incorpora la incertidumbre epistémica mediante términos de desviación estándar en los coeficientes constantes dependientes de la región, representando variaciones regionales en la atenuación anelástica, la escala de $V_{S30}$ y los términos constantes. Para regiones no incluidas directamente en el modelo, se asigna un rango de incertidumbre epistémica que abarque estas variaciones.

%   \item \textbf{Kuehn et al.\ (2020).}  
%   Emplea un enfoque bayesiano para cuantificar formalmente la incertidumbre epistémica a partir de muestras de la distribución posterior de los coeficientes de regresión. Este enfoque permite calcular medianas modeladas, desviaciones estándar aleatorias e incertidumbre epistémica. Además, evalúa la incertidumbre en la magnitud de punto de ruptura, la atenuación anelástica y la amplificación del sitio.

%   \item \textbf{Abrahamson et al.\ (2016) y Abrahamson et al.\ (2018).}  
%   El modelo de 2016 considera la incertidumbre epistémica mediante factores de escala alternativos para la mediana, basados en el rango de las medias de los términos específicos de eventos regionales. También ajusta el punto de ruptura en el escalado de magnitud en $\pm 0.5$ unidades. En 2018, se emplea la base de datos NGA-Subduction, heredando las consideraciones de 2016 y reflejando la incertidumbre en los coeficientes de regresión y factores de escala alternativos.
% \end{itemize}

% \subsubsection{Variabilidad aleatoria}
% La variabilidad aleatoria en los GMM para zonas de subducción se aborda mediante distintos parámetros estadísticos y metodologías:

% \begin{itemize}
%   \item \textbf{Montalva et al.\ (2017).}  
%   Evalúa la sensibilidad de los coeficientes del modelo al tamaño del conjunto de datos mediante un análisis de remuestreo, generando distribuciones de coeficientes a partir de regresiones repetidas. Además, compara las desviaciones estándar del modelo con las de otras regiones y entornos tectónicos, desglosando componentes como la desviación estándar entre eventos ($\tau$), dentro del evento ($\phi$), dentro del evento de una sola estación ($\phi_{SS}$) y de sitio a sitio ($\phi_{S2S}$).

%   \item \textbf{Parker et al.\ (2020).}  
%   Desarrolla modelos con cuatro componentes principales de variabilidad ($\tau$, $\phi$, $\phi_{SS}$ y $\phi_{S2S}$), diferenciando entre eventos de interplaca e intraplaca y diversas medidas de intensidad (\textit{IM}). La variabilidad dentro del evento depende de $V_{S30}$ a distancias menores de 200\,km y de los efectos de trayectoria a distancias mayores.

%   \item \textbf{Kuehn et al.\ (2020).}  
%   Emplea un enfoque bayesiano robusto, utilizando una distribución $t$ de Student para modelar la variabilidad dentro del evento y lograr mayor resistencia frente a valores atípicos. Regionaliza el parámetro $\nu$ para capturar variaciones locales y considera la profundidad de la cuenca en la variabilidad aleatoria, ajustando $\phi$ para evitar contabilizar dos veces ciertos efectos al incluir predictores como $Z_{1.0}$ o $Z_{2.5}$.

%   \item \textbf{Abrahamson et al.\ (2018).}  
%   Incluye componentes de variabilidad entre eventos y dentro del evento, basados en residuos en las regresiones. El modelo refleja adecuadamente los cambios en la atenuación entre sitios de arco frontal y posterior, al no mostrar tendencias significativas con la distancia en los residuos dentro del evento.
% \end{itemize}

% \newpage




% \newpage

% %\begin{landscape}
    



% %\end{landscape}




% \subsection{FSR}
% \label{sec:fsr}



% Los GMPE de NGA-West2 comparten diversos elementos en común. En primer lugar, todos se basan en la base de datos NGA-West2, que contiene más de 21,000 registros de 599 terremotos corticales con magnitudes entre 3.0 y 7.9 \shortcite{Ancheta2014, Bozorgnia2014, Gregor2014}. Esta base de datos incluye eventos ocurridos en China, Japón, el Mediterráneo, Taiwán y el oeste de América del Norte \shortcite{Bozorgnia2014}.

% En cuanto a la medida de intensidad, los cinco modelos emplean \textit{RotD50} para la componente horizontal del movimiento del suelo \shortcite{Gregor2014, Ancheta2014, Bozorgnia2014}. \textit{RotD50} corresponde al percentil 50 del componente horizontal combinado, independiente de la orientación y dependiente del período \shortcite{Ancheta2014}. Los modelos predicen la aceleración máxima del suelo (PGA), la velocidad máxima del suelo (PGV) y la aceleración espectral de respuesta pseudo-absoluta (PSA) para, al menos, 21 períodos de oscilador que abarcan desde 0.01 hasta 10 segundos \shortcite{Bozorgnia2014}.

% Uno de los aspectos considerados es el escalamiento no lineal de la magnitud, al reconocer que la tasa de aumento del movimiento del suelo con la magnitud no es constante. Asimismo, se contempla que la tasa de disminución del movimiento del suelo con la distancia varía en función de la magnitud del evento \shortcite{Bozorgnia2014}.

% Todos los modelos incluyen términos para la amplificación del sitio. Emplean $V_{S30}$, que es la velocidad promedio de la onda de corte en los 30\,m superiores, como parámetro del sitio \shortcite{Bozorgnia2014, Gregor2014}. Algunos modelos también contemplan otros parámetros, como la profundidad al límite de la onda de corte de 1.0\,km/s (\(Z_{1.0}\)) \shortcite{Gregor2014}.

% Cuatro de los cinco modelos (ASK, BSSA, CB y CY) incluyen ajustes regionales relacionados con la atenuación o la respuesta del sitio \shortcite{Bozorgnia2014, Ancheta2014}. Dichos ajustes se basan en el análisis de subconjuntos de datos regionales dentro de la base de datos NGA-West2.

% En la Tabla \ref{tab: Variables usadas NGA-West2} se presenta una comparación entre los parámetros de entrada utilizados por cada uno de los GMPE. En la Figura \ref{fig:Distancias usadas} se ilustran las distancias aplicadas para los cálculos asociados a las GMPE.

% \subsubsection{Comparación y diferencias entre los GMPE de NGA-West2}

% Los cinco modelos de movimiento del suelo (GMPE) del proyecto NGA-West2 exhiben diferencias en varios aspectos de su formulación:

% \begin{itemize}
%     \item \textbf{Métricas de distancia:} Mientras que ASK, CB, CY e IM emplean tanto la distancia más cercana al plano de ruptura (RRUP) como la distancia más cercana a la proyección horizontal de dicho plano (RJB), BSSA se basa exclusivamente en RJB \shortcite{Gregor2014}.

%     \item \textbf{Caracterización de la amplificación no lineal del sitio:} ASK y CY utilizan la aceleración espectral en un sitio de referencia para modelar la no linealidad del suelo, mientras que BSSA y CB emplean la aceleración máxima del suelo (PGA) con el mismo propósito \shortcite{Gregor2014}.

%     \item \textbf{Parámetros de atenuación:} ASK, CB y CY introducen ajustes en la atenuación a grandes distancias, modificando los términos que controlan la reducción del movimiento del suelo en función de la distancia. BSSA incluye ajustes de atenuación regional \shortcite{Gregor2014}.

%     \item \textbf{Términos de respuesta del sitio:} Los modelos difieren también en el tratamiento de los términos de respuesta del sitio, que influyen en la amplificación del movimiento del suelo de acuerdo con las condiciones locales. ASK, BSSA, CB y CY contemplan ajustes para distintas regiones \shortcite{Gregor2014}.
% \end{itemize}

% En la Tabla \ref{tabla: Rango de aplicabilidad NGA-West2} se detallan las condiciones bajo las cuales los modelos de NGA-West2 son aplicables para fallas inversas.





% \begin{figure}
%     \centering
%         % Primera imagen
%     \begin{minipage}{0.48\textwidth}
%         \centering
%         \includegraphics[width=\textwidth]{figuras_tesis/contexto_geologico/Distancias_1.png}
%     \end{minipage}
%     \hfill
%     % Segunda imagen
%     \begin{minipage}{0.48\textwidth}
%         \centering
%         \includegraphics[width=\textwidth]{figuras_tesis/contexto_geologico/Distancias_2.png}
%     \end{minipage}

%     % Ajustar el ancho del caption
%     \captionsetup{width=\textwidth}
%     \caption[Distancias empleadas en sismología]{Esquema ilustrativo de las medidas de distancia utilizadas en un modelo sismológico. Izquierda: esquema de la fuente del terremoto y medidas de distancia empleando una sección vertical a través de un plano de ruptura de falla. La extensión del plano de ruptura de la falla ($L$) se mide siguiendo la dirección del rumbo (perpendicular al plano de la página). Derecha: vista en planta que ilustra la definición de $R_Y$. Recuperado de Kaklamanos et al. (2011).}
%     \label{fig:Distancias usadas}
% \end{figure}








\section{Gráficos extra}

\subsection{Mapas de Amenaza DSHA}

Se muestran en esta sección los resultados de los modelos $M_w$ 7.4, 7.3, 7.2 y 6.8 en sus tres variaciones, para los modelos con ruptura de profundidad mínima de 0 km y profundidad mínia de 5 km.

\subsubsection{FSR \texorpdfstring{$M_w$ 7.4}{Mw 7.4}}



\begin{figure}
    \centering
    \includegraphics[width=0.7\linewidth]{figuras_tesis/resultados_amenaza/DSHA/PGA FSR 74 DIP 34 PROF MIN 0 KM.png}
    \caption{Mapa de amenaza sísmica para PGA, considerando un escenario de ruptura en FSR con $M_w$ 7.4, Dip 34° y profundidad mínima de ruptura de 0 km.}
    \label{fig: Mapa DSHA PGA FSR 74 DIP 34 PROF MIN 0 KM}
\end{figure}

\begin{figure}
    \centering
    \includegraphics[width=0.7\linewidth]{figuras_tesis/resultados_amenaza/DSHA/PGA FSR 74 DIP 60 PROF MIN 0 KM.png}
    \caption{Mapa de amenaza sísmica para PGA, considerando un escenario de ruptura en FSR con $M_w$ 7.4, Dip 60° y profundidad mínima de ruptura de 0 km.}
    \label{fig: Mapa DSHA PGA FSR 74 DIP 60 PROF MIN 0 KM}
\end{figure}

\begin{figure}
    \centering
    \includegraphics[width=0.7\linewidth]{figuras_tesis/resultados_amenaza/DSHA/PGA FSR 74 DIP 34 PROF MIN 5 KM.png}
    \caption{Mapa de amenaza sísmica para PGA, considerando un escenario de ruptura en FSR con $M_w$ 7.4, Dip 34° y profundidad mínima de ruptura de 5 km.}
    \label{fig: Mapa DSHA PGA FSR 74 DIP 34 PROF MIN 5 KM}
\end{figure}

\begin{figure}
    \centering
    \includegraphics[width=0.7\linewidth]{figuras_tesis/resultados_amenaza/DSHA/PGA FSR 74 DIP 60 PROF MIN 5 KM.png}
    \caption{Mapa de amenaza sísmica para PGA, considerando un escenario de ruptura en FSR con $M_w$ 7.4, Dip 60° y profundidad mínima de ruptura de 5 km.}
    \label{fig: Mapa DSHA PGA FSR 74 DIP 60 PROF MIN 5 KM}
\end{figure}










\begin{figure}
    \centering
    \includegraphics[width=0.7\linewidth]{figuras_tesis/resultados_amenaza/DSHA/SA02 FSR 74 DIP 34 PROF MIN 0 KM.png}
    \caption{Mapa de amenaza sísmica para SA(0.2), considerando un escenario de ruptura en FSR con $M_w$ 7.4, Dip 34° y profundidad mínima de ruptura de 0 km.}
    \label{fig: Mapa DSHA SA(0.2) FSR 74 DIP 34 PROF MIN 0 KM}
\end{figure}

\begin{figure}
    \centering
    \includegraphics[width=0.7\linewidth]{figuras_tesis/resultados_amenaza/DSHA/SA02 FSR 74 DIP 60 PROF MIN 0 KM.png}
    \caption{Mapa de amenaza sísmica para SA(0.2), considerando un escenario de ruptura en FSR con $M_w$ 7.4, Dip 60° y profundidad mínima de ruptura de 0 km.}
    \label{fig: Mapa DSHA SA(0.2) FSR 74 DIP 60 PROF MIN 0 KM}
\end{figure}

\begin{figure}
    \centering
    \includegraphics[width=0.7\linewidth]{figuras_tesis/resultados_amenaza/DSHA/SA02 FSR 74 DIP 34 PROF MIN 5 KM.png}
    \caption{Mapa de amenaza sísmica para SA(0.2), considerando un escenario de ruptura en FSR con $M_w$ 7.4, Dip 34° y profundidad mínima de ruptura de 5 km.}
    \label{fig: Mapa DSHA SA(0.2) FSR 74 DIP 34 PROF MIN 5 KM}
\end{figure}

\begin{figure}
    \centering
    \includegraphics[width=0.7\linewidth]{figuras_tesis/resultados_amenaza/DSHA/SA02 FSR 74 DIP 60 PROF MIN 5 KM.png}
    \caption{Mapa de amenaza sísmica para SA(0.2), considerando un escenario de ruptura en FSR con $M_w$ 7.4, Dip 60° y profundidad mínima de ruptura de 5 km.}
    \label{fig: Mapa DSHA SA(0.2) FSR 74 DIP 60 PROF MIN 5 KM}
\end{figure}








\begin{figure}
    \centering
    \includegraphics[width=0.7\linewidth]{figuras_tesis/resultados_amenaza/DSHA/SA10 FSR 74 DIP 34 PROF MIN 0 KM.png}
    \caption{Mapa de amenaza sísmica para SA(1.0), considerando un escenario de ruptura en FSR con $M_w$ 7.4, Dip 34° y profundidad mínima de ruptura de 0 km.}
    \label{fig: Mapa DSHA SA(1.0) FSR 74 DIP 34 PROF MIN 0 KM}
\end{figure}

\begin{figure}
    \centering
    \includegraphics[width=0.7\linewidth]{figuras_tesis/resultados_amenaza/DSHA/SA10 FSR 74 DIP 60 PROF MIN 0 KM.png}
    \caption{Mapa de amenaza sísmica para SA(1.0), considerando un escenario de ruptura en FSR con $M_w$ 7.4, Dip 60° y profundidad mínima de ruptura de 0 km.}
    \label{fig: Mapa DSHA SA(1.0) FSR 74 DIP 60 PROF MIN 0 KM}
\end{figure}

\begin{figure}
    \centering
    \includegraphics[width=0.7\linewidth]{figuras_tesis/resultados_amenaza/DSHA/SA10 FSR 74 DIP 34 PROF MIN 5 KM.png}
    \caption{Mapa de amenaza sísmica para SA(1.0), considerando un escenario de ruptura en FSR con $M_w$ 7.4, Dip 34° y profundidad mínima de ruptura de 5 km.}
    \label{fig: Mapa DSHA SA(1.0) FSR 74 DIP 34 PROF MIN 5 KM}
\end{figure}

\begin{figure}
    \centering
    \includegraphics[width=0.7\linewidth]{figuras_tesis/resultados_amenaza/DSHA/SA10 FSR 74 DIP 60 PROF MIN 5 KM.png}
    \caption{Mapa de amenaza sísmica para SA(1.0), considerando un escenario de ruptura en FSR con $M_w$ 7.4, Dip 60° y profundidad mínima de ruptura de 5 km.}
    \label{fig: Mapa DSHA SA(1.0) FSR 74 DIP 60 PROF MIN 5 KM}
\end{figure}


















\subsubsection{FSR \texorpdfstring{$M_w$ 7.3}{Mw 7.3}}



\begin{figure}
    \centering
    \includegraphics[width=0.7\linewidth]{figuras_tesis/resultados_amenaza/DSHA/PGA FSR 73 DIP 34 PROF MIN 0 KM.png}
    \caption{Mapa de amenaza sísmica para PGA, considerando un escenario de ruptura en FSR con $M_w$ 7.3, Dip 34° y profundidad mínima de ruptura de 0 km.}
    \label{fig: Mapa DSHA PGA FSR 73 DIP 34 PROF MIN 0 KM}
\end{figure}

\begin{figure}
    \centering
    \includegraphics[width=0.7\linewidth]{figuras_tesis/resultados_amenaza/DSHA/PGA FSR 73 DIP 60 PROF MIN 0 KM.png}
    \caption{Mapa de amenaza sísmica para PGA, considerando un escenario de ruptura en FSR con $M_w$ 7.3, Dip 60° y profundidad mínima de ruptura de 0 km.}
    \label{fig: Mapa DSHA PGA FSR 73 DIP 60 PROF MIN 0 KM}
\end{figure}

\begin{figure}
    \centering
    \includegraphics[width=0.7\linewidth]{figuras_tesis/resultados_amenaza/DSHA/PGA FSR 73 DIP 34 PROF MIN 5 KM.png}
    \caption{Mapa de amenaza sísmica para PGA, considerando un escenario de ruptura en FSR con $M_w$ 7.3, Dip 34° y profundidad mínima de ruptura de 5 km.}
    \label{fig: Mapa DSHA PGA FSR 73 DIP 34 PROF MIN 5 KM}
\end{figure}

\begin{figure}
    \centering
    \includegraphics[width=0.7\linewidth]{figuras_tesis/resultados_amenaza/DSHA/PGA FSR 73 DIP 60 PROF MIN 5 KM.png}
    \caption{Mapa de amenaza sísmica para PGA, considerando un escenario de ruptura en FSR con $M_w$ 7.3, Dip 34° y profundidad mínima de ruptura de 0 km.}
    \label{fig: Mapa DSHA PGA FSR 73 DIP 60 PROF MIN 5 KM}
\end{figure}










\begin{figure}
    \centering
    \includegraphics[width=0.7\linewidth]{figuras_tesis/resultados_amenaza/DSHA/SA02 FSR 73 DIP 34 PROF MIN 0 KM.png}
    \caption{Mapa de amenaza sísmica para SA(0.2), considerando un escenario de ruptura en FSR con $M_w$ 7.3, Dip 34° y profundidad mínima de ruptura de 0 km.}
    \label{fig: Mapa DSHA SA(0.2) FSR 73 DIP 34 PROF MIN 0 KM}
\end{figure}

\begin{figure}
    \centering
    \includegraphics[width=0.7\linewidth]{figuras_tesis/resultados_amenaza/DSHA/SA02 FSR 73 DIP 60 PROF MIN 0 KM.png}
    \caption{Mapa de amenaza sísmica para SA(0.2), considerando un escenario de ruptura en FSR con $M_w$ 7.3, Dip 60° y profundidad mínima de ruptura de 0 km.}
    \label{fig: Mapa DSHA SA(0.2) FSR 73 DIP 60 PROF MIN 0 KM}
\end{figure}

\begin{figure}
    \centering
    \includegraphics[width=0.7\linewidth]{figuras_tesis/resultados_amenaza/DSHA/SA02 FSR 73 DIP 34 PROF MIN 5 KM.png}
    \caption{Mapa de amenaza sísmica para SA(0.2), considerando un escenario de ruptura en FSR con $M_w$ 7.3, Dip 34° y profundidad mínima de ruptura de 5 km.}
    \label{fig: Mapa DSHA SA(0.2) FSR 73 DIP 34 PROF MIN 5 KM}
\end{figure}

\begin{figure}
    \centering
    \includegraphics[width=0.7\linewidth]{figuras_tesis/resultados_amenaza/DSHA/SA02 FSR 73 DIP 60 PROF MIN 5 KM.png}
    \caption{Mapa de amenaza sísmica para SA(0.2), considerando un escenario de ruptura en FSR con $M_w$ 7.3, Dip 60° y profundidad mínima de ruptura de 5 km.}
    \label{fig: Mapa DSHA SA(0.2) FSR 73 DIP 60 PROF MIN 5 KM}
\end{figure}








\begin{figure}
    \centering
    \includegraphics[width=0.7\linewidth]{figuras_tesis/resultados_amenaza/DSHA/SA10 FSR 73 DIP 34 PROF MIN 0 KM.png}
    \caption{Mapa de amenaza sísmica para SA(1.0), considerando un escenario de ruptura en FSR con $M_w$ 7.3, Dip 34° y profundidad mínima de ruptura de 0 km.}
    \label{fig: Mapa DSHA SA(1.0) FSR 73 DIP 34 PROF MIN 0 KM}
\end{figure}

\begin{figure}
    \centering
    \includegraphics[width=0.7\linewidth]{figuras_tesis/resultados_amenaza/DSHA/SA10 FSR 73 DIP 60 PROF MIN 0 KM.png}
    \caption{Mapa de amenaza sísmica para SA(1.0), considerando un escenario de ruptura en FSR con $M_w$ 7.3, Dip 60° y profundidad mínima de ruptura de 0 km.}
    \label{fig: Mapa DSHA SA(1.0) FSR 73 DIP 60 PROF MIN 0 KM}
\end{figure}

\begin{figure}
    \centering
    \includegraphics[width=0.7\linewidth]{figuras_tesis/resultados_amenaza/DSHA/SA10 FSR 73 DIP 34 PROF MIN 5 KM.png}
    \caption{Mapa de amenaza sísmica para SA(1.0), considerando un escenario de ruptura en FSR con $M_w$ 7.3, Dip 34° y profundidad mínima de ruptura de 5 km.}
    \label{fig: Mapa DSHA SA(1.0) FSR 73 DIP 34 PROF MIN 5 KM}
\end{figure}

\begin{figure}
    \centering
    \includegraphics[width=0.7\linewidth]{figuras_tesis/resultados_amenaza/DSHA/SA10 FSR 73 DIP 60 PROF MIN 5 KM.png}
    \caption{Mapa de amenaza sísmica para SA(1.0), considerando un escenario de ruptura en FSR con $M_w$ 7.3, Dip 60° y profundidad mínima de ruptura de 5 km.}
    \label{fig: Mapa DSHA SA(1.0) FSR 73 DIP 60 PROF MIN 5 KM}
\end{figure}
















\subsubsection{FSR \texorpdfstring{$M_w$ 7.2}{Mw 7.2}}



\begin{figure}
    \centering
    \includegraphics[width=0.7\linewidth]{figuras_tesis/resultados_amenaza/DSHA/PGA FSR 72 DIP 34 PROF MIN 0 KM.png}
    \caption{Mapa de amenaza sísmica para PGA, considerando un escenario de ruptura en FSR con $M_w$ 7.2, Dip 34° y profundidad mínima de ruptura de 0 km.}
    \label{fig: Mapa DSHA PGA FSR 72 DIP 34 PROF MIN 0 KM}
\end{figure}

\begin{figure}
    \centering
    \includegraphics[width=0.7\linewidth]{figuras_tesis/resultados_amenaza/DSHA/PGA FSR 72 DIP 60 PROF MIN 0 KM.png}
    \caption{Mapa de amenaza sísmica para PGA, considerando un escenario de ruptura en FSR con $M_w$ 7.2, Dip 60° y profundidad mínima de ruptura de 0 km.}
    \label{fig: Mapa DSHA PGA FSR 72 DIP 60 PROF MIN 0 KM}
\end{figure}

\begin{figure}
    \centering
    \includegraphics[width=0.7\linewidth]{figuras_tesis/resultados_amenaza/DSHA/PGA FSR 72 DIP 34 PROF MIN 5 KM.png}
    \caption{Mapa de amenaza sísmica para PGA, considerando un escenario de ruptura en FSR con $M_w$ 7.2, Dip 34° y profundidad mínima de ruptura de 5 km.}
    \label{fig: Mapa DSHA PGA FSR 72 DIP 34 PROF MIN 5 KM}
\end{figure}

\begin{figure}
    \centering
    \includegraphics[width=0.7\linewidth]{figuras_tesis/resultados_amenaza/DSHA/PGA FSR 72 DIP 60 PROF MIN 5 KM.png}
    \caption{Mapa de amenaza sísmica para PGA, considerando un escenario de ruptura en FSR con $M_w$ 7.2, Dip 60° y profundidad mínima de ruptura de 5 km.}
    \label{fig: Mapa DSHA PGA FSR 72 DIP 60 PROF MIN 5 KM}
\end{figure}










\begin{figure}
    \centering
    \includegraphics[width=0.7\linewidth]{figuras_tesis/resultados_amenaza/DSHA/SA02 FSR 72 DIP 34 PROF MIN 0 KM.png}
    \caption{Mapa de amenaza sísmica para SA(0.2), considerando un escenario de ruptura en FSR con $M_w$ 7.2, Dip 34° y profundidad mínima de ruptura de 0 km.}
    \label{fig: Mapa DSHA SA(0.2) FSR 72 DIP 34 PROF MIN 0 KM}
\end{figure}

\begin{figure}
    \centering
    \includegraphics[width=0.7\linewidth]{figuras_tesis/resultados_amenaza/DSHA/SA02 FSR 72 DIP 60 PROF MIN 0 KM.png}
    \caption{Mapa de amenaza sísmica para SA(0.2), considerando un escenario de ruptura en FSR con $M_w$ 7.2, Dip 60° y profundidad mínima de ruptura de 0 km.}
    \label{fig: Mapa DSHA SA(0.2) FSR 72 DIP 60 PROF MIN 0 KM}
\end{figure}

\begin{figure}
    \centering
    \includegraphics[width=0.7\linewidth]{figuras_tesis/resultados_amenaza/DSHA/SA02 FSR 72 DIP 34 PROF MIN 5 KM.png}
    \caption{Mapa de amenaza sísmica para SA(0.2), considerando un escenario de ruptura en FSR con $M_w$ 7.2, Dip 34° y profundidad mínima de ruptura de 5 km.}
    \label{fig: Mapa DSHA SA(0.2) FSR 72 DIP 34 PROF MIN 5 KM}
\end{figure}

\begin{figure}
    \centering
    \includegraphics[width=0.7\linewidth]{figuras_tesis/resultados_amenaza/DSHA/SA02 FSR 72 DIP 60 PROF MIN 5 KM.png}
    \caption{Mapa de amenaza sísmica para SA(0.2), considerando un escenario de ruptura en FSR con $M_w$ 7.2, Dip 60° y profundidad mínima de ruptura de 5 km.}
    \label{fig: Mapa DSHA SA(0.2) FSR 72 DIP 60 PROF MIN 5 KM}
\end{figure}








\begin{figure}
    \centering
    \includegraphics[width=0.7\linewidth]{figuras_tesis/resultados_amenaza/DSHA/SA10 FSR 72 DIP 34 PROF MIN 0 KM.png}
    \caption{Mapa de amenaza sísmica para SA(1.0), considerando un escenario de ruptura en FSR con $M_w$ 7.2, Dip 34° y profundidad mínima de ruptura de 0 km.}
    \label{fig: Mapa DSHA SA(1.0) FSR 72 DIP 34 PROF MIN 0 KM}
\end{figure}

\begin{figure}
    \centering
    \includegraphics[width=0.7\linewidth]{figuras_tesis/resultados_amenaza/DSHA/SA10 FSR 72 DIP 60 PROF MIN 0 KM.png}
    \caption{Mapa de amenaza sísmica para SA(1.0), considerando un escenario de ruptura en FSR con $M_w$ 7.2, Dip 60° y profundidad mínima de ruptura de 0 km.}
    \label{fig: Mapa DSHA SA(1.0) FSR 72 DIP 60 PROF MIN 0 KM}
\end{figure}

\begin{figure}
    \centering
    \includegraphics[width=0.7\linewidth]{figuras_tesis/resultados_amenaza/DSHA/SA10 FSR 72 DIP 34 PROF MIN 5 KM.png}
    \caption{Mapa de amenaza sísmica para SA(1.0), considerando un escenario de ruptura en FSR con $M_w$ 7.2, Dip 34° y profundidad mínima de ruptura de 5 km.}
    \label{fig: Mapa DSHA SA(1.0) FSR 72 DIP 34 PROF MIN 5 KM}
\end{figure}

\begin{figure}
    \centering
    \includegraphics[width=0.7\linewidth]{figuras_tesis/resultados_amenaza/DSHA/SA10 FSR 72 DIP 60 PROF MIN 5 KM.png}
    \caption{Mapa de amenaza sísmica para SA(1.0), considerando un escenario de ruptura en FSR con $M_w$ 7.2, Dip 60° y profundidad mínima de ruptura de 5 km.}
    \label{fig: Mapa DSHA SA(1.0) FSR 72 DIP 60 PROF MIN 5 KM}
\end{figure}

























\subsubsection{FSR \texorpdfstring{$M_w$ 6.8 Caso 1}{Mw 6.8 caso 1}}

\begin{figure}
    \centering
    \includegraphics[width=0.7\linewidth]{figuras_tesis/resultados_amenaza/DSHA/PGA FSR 68 DIP 34 CASO 1 PROF MIN 0 KM.png}
    \caption{Mapa de amenaza sísmica para PGA, considerando un escenario de ruptura en FSR con $M_w$ 6.8, Dip 34°, caso 1 y profundidad mínima de ruptura de 0 km.}
    \label{fig: Mapa DSHA PGA FSR 68 CASO 1 DIP 34 PROF MIN 0 KM}
\end{figure}

\begin{figure}
    \centering
    \includegraphics[width=0.7\linewidth]{figuras_tesis/resultados_amenaza/DSHA/PGA FSR 68 DIP 60 CASO 1 PROF MIN 0 KM.png}
    \caption{Mapa de amenaza sísmica para PGA, considerando un escenario de ruptura en FSR con $M_w$ 6.8, Dip 60°, caso 1 y profundidad mínima de ruptura de 0 km.}
    \label{fig: Mapa DSHA PGA FSR 68 CASO 1 DIP 60 PROF MIN 0 KM}
\end{figure}

\begin{figure}
    \centering
    \includegraphics[width=0.7\linewidth]{figuras_tesis/resultados_amenaza/DSHA/PGA FSR 68 DIP 34 CASO 1 PROF MIN 5 KM.png}
    \caption{Mapa de amenaza sísmica para PGA, considerando un escenario de ruptura en FSR con $M_w$ 6.8, Dip 34°, caso 1 y profundidad mínima de ruptura de 5 km.}
    \label{fig: Mapa DSHA PGA FSR 68 CASO 1 DIP 34 PROF MIN 5 KM}
\end{figure}

\begin{figure}
    \centering
    \includegraphics[width=0.7\linewidth]{figuras_tesis/resultados_amenaza/DSHA/PGA FSR 68 DIP 60 CASO 1 PROF MIN 5 KM.png}
    \caption{Mapa de amenaza sísmica para PGA, considerando un escenario de ruptura en FSR con $M_w$ 6.8, Dip 60°, caso 1 y profundidad mínima de ruptura de 5 km.}
    \label{fig: Mapa DSHA PGA FSR 68 CASO 1 DIP 60 PROF MIN 5 KM}
\end{figure}









\begin{figure}
    \centering
    \includegraphics[width=0.7\linewidth]{figuras_tesis/resultados_amenaza/DSHA/SA02 FSR 68 DIP 34 CASO 1 PROF MIN 0 KM.png}
    \caption{Mapa de amenaza sísmica para SA(0.2), considerando un escenario de ruptura en FSR con $M_w$ 6.8, Dip 34°, caso 1 y profundidad mínima de ruptura de 0 km.}
    \label{fig: Mapa DSHA SA(0.2) FSR 68 CASO 1 DIP 34 PROF MIN 0 KM}
\end{figure}

\begin{figure}
    \centering
    \includegraphics[width=0.7\linewidth]{figuras_tesis/resultados_amenaza/DSHA/SA02 FSR 68 DIP 60 CASO 1 PROF MIN 0 KM.png}
    \caption{Mapa de amenaza sísmica para SA(0.2), considerando un escenario de ruptura en FSR con $M_w$ 6.8, Dip 60°, caso 1 y profundidad mínima de ruptura de 0 km.}
    \label{fig: Mapa DSHA SA(0.2) FSR 68 CASO 1 DIP 60 PROF MIN 0 KM}
\end{figure}

\begin{figure}
    \centering
    \includegraphics[width=0.7\linewidth]{figuras_tesis/resultados_amenaza/DSHA/SA02 FSR 68 DIP 34 CASO 1 PROF MIN 5 KM.png}
    \caption{Mapa de amenaza sísmica para SA(0.2), considerando un escenario de ruptura en FSR con $M_w$ 6.8, Dip 34°, caso 1 y profundidad mínima de ruptura de 5 km.}
    \label{fig: Mapa DSHA SA(0.2) FSR 68 CASO 1 DIP 34 PROF MIN 5 KM}
\end{figure}

\begin{figure}
    \centering
    \includegraphics[width=0.7\linewidth]{figuras_tesis/resultados_amenaza/DSHA/SA02 FSR 68 DIP 60 CASO 1 PROF MIN 5 KM.png}
    \caption{Mapa de amenaza sísmica para SA(0.2), considerando un escenario de ruptura en FSR con $M_w$ 6.8, Dip 60°, caso 1 y profundidad mínima de ruptura de 5 km.}
    \label{fig: Mapa DSHA SA(0.2) FSR 68 CASO 1 DIP 60 PROF MIN 5 KM}
\end{figure}









\begin{figure}
    \centering
    \includegraphics[width=0.7\linewidth]{figuras_tesis/resultados_amenaza/DSHA/SA10 FSR 68 DIP 34 CASO 1 PROF MIN 0 KM.png}
    \caption{Mapa de amenaza sísmica para SA(1.0), considerando un escenario de ruptura en FSR con $M_w$ 6.8, Dip 34°, caso 1 y profundidad mínima de ruptura de 0 km.}
    \label{fig: Mapa DSHA SA(1.0) FSR 68 CASO 1 DIP 34 PROF MIN 0 KM}
\end{figure}

\begin{figure}
    \centering
    \includegraphics[width=0.7\linewidth]{figuras_tesis/resultados_amenaza/DSHA/SA10 FSR 68 DIP 60 CASO 1 PROF MIN 0 KM.png}
    \caption{Mapa de amenaza sísmica para SA(1.0), considerando un escenario de ruptura en FSR con $M_w$ 6.8, Dip 60°, caso 1 y profundidad mínima de ruptura de 0 km.}
    \label{fig: Mapa DSHA SA(1.0) FSR 68 CASO 1 DIP 60 PROF MIN 0 KM}
\end{figure}

\begin{figure}
    \centering
    \includegraphics[width=0.7\linewidth]{figuras_tesis/resultados_amenaza/DSHA/SA10 FSR 68 DIP 34 CASO 1 PROF MIN 5 KM.png}
    \caption{Mapa de amenaza sísmica para SA(1.0), considerando un escenario de ruptura en FSR con $M_w$ 6.8, Dip 34°, caso 1 y profundidad mínima de ruptura de 5 km.}
    \label{fig: Mapa DSHA SA(1.0) FSR 68 CASO 1 DIP 34 PROF MIN 5 KM}
\end{figure}

\begin{figure}
    \centering
    \includegraphics[width=0.7\linewidth]{figuras_tesis/resultados_amenaza/DSHA/SA10 FSR 68 DIP 60 CASO 1 PROF MIN 5 KM.png}
    \caption{Mapa de amenaza sísmica para SA(1.0), considerando un escenario de ruptura en FSR con $M_w$ 6.8, Dip 60°, caso 1 y profundidad mínima de ruptura de 5 km.}
    \label{fig: Mapa DSHA SA(1.0) FSR 68 CASO 1 DIP 60 PROF MIN 5 KM}
\end{figure}



\subsubsection{FSR \texorpdfstring{$M_w$ 6.8 Caso 2}{Mw 6.8 caso 2}}

\begin{figure}
    \centering
    \includegraphics[width=0.7\linewidth]{figuras_tesis/resultados_amenaza/DSHA/PGA FSR 68 DIP 34 CASO 2 PROF MIN 0 KM.png}
    \caption{Mapa de amenaza sísmica para PGA, considerando un escenario de ruptura en FSR con $M_w$ 6.8, Dip 34°, caso 2 y profundidad mínima de ruptura de 0 km.}
    \label{fig: Mapa DSHA PGA FSR 68 CASO 2 DIP 34 PROF MIN 0 KM}
\end{figure}

\begin{figure}
    \centering
    \includegraphics[width=0.7\linewidth]{figuras_tesis/resultados_amenaza/DSHA/PGA FSR 68 DIP 60 CASO 2 PROF MIN 0 KM.png}
    \caption{Mapa de amenaza sísmica para PGA, considerando un escenario de ruptura en FSR con $M_w$ 6.8, Dip 60°, caso 2 y profundidad mínima de ruptura de 0 km.}
    \label{fig: Mapa DSHA PGA FSR 68 CASO 2 DIP 60 PROF MIN 0 KM}
\end{figure}

\begin{figure}
    \centering
    \includegraphics[width=0.7\linewidth]{figuras_tesis/resultados_amenaza/DSHA/PGA FSR 68 DIP 34 CASO 2 PROF MIN 5 KM.png}
    \caption{Mapa de amenaza sísmica para PGA, considerando un escenario de ruptura en FSR con $M_w$ 6.8, Dip 34°, caso 2 y profundidad mínima de ruptura de 5 km.}
    \label{fig: Mapa DSHA PGA FSR 68 CASO 2 DIP 34 PROF MIN 5 KM}
\end{figure}

\begin{figure}
    \centering
    \includegraphics[width=0.7\linewidth]{figuras_tesis/resultados_amenaza/DSHA/PGA FSR 68 DIP 60 CASO 2 PROF MIN 5 KM.png}
    \caption{Mapa de amenaza sísmica para PGA, considerando un escenario de ruptura en FSR con $M_w$ 6.8, Dip 60°, caso 2 y profundidad mínima de ruptura de 5 km.}
    \label{fig: Mapa DSHA PGA FSR 68 CASO 2 DIP 60 PROF MIN 5 KM}
\end{figure}







\begin{figure}
    \centering
    \includegraphics[width=0.7\linewidth]{figuras_tesis/resultados_amenaza/DSHA/SA02 FSR 68 DIP 34 CASO 2 PROF MIN 0 KM.png}
    \caption{Mapa de amenaza sísmica para SA(0.2), considerando un escenario de ruptura en FSR con $M_w$ 6.8, Dip 34°, caso 2 y profundidad mínima de ruptura de 0 km.}
    \label{fig: Mapa DSHA SA(0.2) FSR 68 CASO 2 DIP 34 PROF MIN 0 KM}
\end{figure}

\begin{figure}
    \centering
    \includegraphics[width=0.7\linewidth]{figuras_tesis/resultados_amenaza/DSHA/SA02 FSR 68 DIP 60 CASO 2 PROF MIN 0 KM.png}
    \caption{Mapa de amenaza sísmica para SA(0.2), considerando un escenario de ruptura en FSR con $M_w$ 6.8, Dip 60°, caso 2 y profundidad mínima de ruptura de 0 km.}
    \label{fig: Mapa DSHA SA(0.2) FSR 68 CASO 2 DIP 60 PROF MIN 0 KM}
\end{figure}

\begin{figure}
    \centering
    \includegraphics[width=0.7\linewidth]{figuras_tesis/resultados_amenaza/DSHA/SA02 FSR 68 DIP 34 CASO 2 PROF MIN 5 KM.png}
    \caption{Mapa de amenaza sísmica para SA(0.2), considerando un escenario de ruptura en FSR con $M_w$ 6.8, Dip 34°, caso 2 y profundidad mínima de ruptura de 5 km.}
    \label{fig: Mapa DSHA SA(0.2) FSR 68 CASO 2 DIP 34 PROF MIN 5 KM}
\end{figure}

\begin{figure}
    \centering
    \includegraphics[width=0.7\linewidth]{figuras_tesis/resultados_amenaza/DSHA/SA02 FSR 68 DIP 60 CASO 2 PROF MIN 5 KM.png}
    \caption{Mapa de amenaza sísmica para SA(0.2), considerando un escenario de ruptura en FSR con $M_w$ 6.8, Dip 60°, caso 2 y profundidad mínima de ruptura de 5 km.}
    \label{fig: Mapa DSHA SA(0.2) FSR 68 CASO 2 DIP 60 PROF MIN 5 KM}
\end{figure}







\begin{figure}
    \centering
    \includegraphics[width=0.7\linewidth]{figuras_tesis/resultados_amenaza/DSHA/SA10 FSR 68 DIP 34 CASO 2 PROF MIN 0 KM.png}
    \caption{Mapa de amenaza sísmica para SA(1.0), considerando un escenario de ruptura en FSR con $M_w$ 6.8, Dip 34°, caso 2 y profundidad mínima de ruptura de 0 km.}
    \label{fig: Mapa DSHA SA(1.0) FSR 68 CASO 2 DIP 34 PROF MIN 0 KM}
\end{figure}

\begin{figure}
    \centering
    \includegraphics[width=0.7\linewidth]{figuras_tesis/resultados_amenaza/DSHA/SA10 FSR 68 DIP 60 CASO 2 PROF MIN 0 KM.png}
    \caption{Mapa de amenaza sísmica para SA(1.0), considerando un escenario de ruptura en FSR con $M_w$ 6.8, Dip 60°, caso 2 y profundidad mínima de ruptura de 0 km.}
    \label{fig: Mapa DSHA SA(1.0) FSR 68 CASO 2 DIP 60 PROF MIN 0 KM}
\end{figure}

\begin{figure}
    \centering
    \includegraphics[width=0.7\linewidth]{figuras_tesis/resultados_amenaza/DSHA/SA10 FSR 68 DIP 34 CASO 2 PROF MIN 5 KM.png}
    \caption{Mapa de amenaza sísmica para SA(1.0), considerando un escenario de ruptura en FSR con $M_w$ 6.8, Dip 34°, caso 2 y profundidad mínima de ruptura de 5 km.}
    \label{fig: Mapa DSHA SA(1.0) FSR 68 CASO 2 DIP 34 PROF MIN 5 KM}
\end{figure}

\begin{figure}
    \centering
    \includegraphics[width=0.7\linewidth]{figuras_tesis/resultados_amenaza/DSHA/SA10 FSR 68 DIP 60 CASO 2 PROF MIN 5 KM.png}
    \caption{Mapa de amenaza sísmica para SA(1.0), considerando un escenario de ruptura en FSR con $M_w$ 6.8, Dip 60°, caso 2 y profundidad mínima de ruptura de 5 km.}
    \label{fig: Mapa DSHA SA(1.0) FSR 68 CASO 2 DIP 60 PROF MIN 5 KM}
\end{figure}









\subsubsection{FSR \texorpdfstring{$M_w$ 6.8 Caso 3}{Mw 6.8 caso 3}}

\begin{figure}
    \centering
    \includegraphics[width=0.7\linewidth]{figuras_tesis/resultados_amenaza/DSHA/PGA FSR 68 DIP 34 CASO 3 PROF MIN 0 KM.png}
    \caption{Mapa de amenaza sísmica para PGA, considerando un escenario de ruptura en FSR con $M_w$ 6.8, Dip 34°, caso 3 y profundidad mínima de ruptura de 0 km.}
    \label{fig: Mapa DSHA PGA FSR 68 CASO 3 DIP 34 PROF MIN 0 KM}
\end{figure}

\begin{figure}
    \centering
    \includegraphics[width=0.7\linewidth]{figuras_tesis/resultados_amenaza/DSHA/PGA FSR 68 DIP 60 CASO 3 PROF MIN 0 KM.png}
    \caption{Mapa de amenaza sísmica para PGA, considerando un escenario de ruptura en FSR con $M_w$ 6.8, Dip 60°, caso 3 y profundidad mínima de ruptura de 0 km.}
    \label{fig: Mapa DSHA PGA FSR 68 CASO 3 DIP 60 PROF MIN 0 KM}
\end{figure}

\begin{figure}
    \centering
    \includegraphics[width=0.7\linewidth]{figuras_tesis/resultados_amenaza/DSHA/PGA FSR 68 DIP 34 CASO 3 PROF MIN 5 KM.png}
    \caption{Mapa de amenaza sísmica para PGA, considerando un escenario de ruptura en FSR con $M_w$ 6.8, Dip 34°, caso 3 y profundidad mínima de ruptura de 5 km.}
    \label{fig: Mapa DSHA PGA FSR 68 CASO 3 DIP 34 PROF MIN 5 KM}
\end{figure}

\begin{figure}
    \centering
    \includegraphics[width=0.7\linewidth]{figuras_tesis/resultados_amenaza/DSHA/PGA FSR 68 DIP 60 CASO 3 PROF MIN 5 KM.png}
    \caption{Mapa de amenaza sísmica para PGA, considerando un escenario de ruptura en FSR con $M_w$ 6.8, Dip 60°, caso 3 y profundidad mínima de ruptura de 5 km.}
    \label{fig: Mapa DSHA PGA FSR 68 CASO 3 DIP 60 PROF MIN 5 KM}
\end{figure}








\begin{figure}
    \centering
    \includegraphics[width=0.7\linewidth]{figuras_tesis/resultados_amenaza/DSHA/SA02 FSR 68 DIP 34 CASO 3 PROF MIN 0 KM.png}
    \caption{Mapa de amenaza sísmica para SA(0.2), considerando un escenario de ruptura en FSR con $M_w$ 6.8, Dip 34°, caso 3 y profundidad mínima de ruptura de 0 km.}
    \label{fig: Mapa DSHA SA(0.2) FSR 68 CASO 3 DIP 34 PROF MIN 0 KM}
\end{figure}

\begin{figure}
    \centering
    \includegraphics[width=0.7\linewidth]{figuras_tesis/resultados_amenaza/DSHA/SA02 FSR 68 DIP 60 CASO 3 PROF MIN 0 KM.png}
    \caption{Mapa de amenaza sísmica para SA(0.2), considerando un escenario de ruptura en FSR con $M_w$ 6.8, Dip 60°, caso 3 y profundidad mínima de ruptura de 0 km.}
    \label{fig: Mapa DSHA SA(0.2) FSR 68 CASO 3 DIP 60 PROF MIN 0 KM}
\end{figure}

\begin{figure}
    \centering
    \includegraphics[width=0.7\linewidth]{figuras_tesis/resultados_amenaza/DSHA/SA02 FSR 68 DIP 34 CASO 3 PROF MIN 5 KM.png}
    \caption{Mapa de amenaza sísmica para SA(0.2), considerando un escenario de ruptura en FSR con $M_w$ 6.8, Dip 34°, caso 3 y profundidad mínima de ruptura de 5 km.}
    \label{fig: Mapa DSHA SA(0.2) FSR 68 CASO 3 DIP 34 PROF MIN 5 KM}
\end{figure}

\begin{figure}
    \centering
    \includegraphics[width=0.7\linewidth]{figuras_tesis/resultados_amenaza/DSHA/SA02 FSR 68 DIP 60 CASO 3 PROF MIN 5 KM.png}
    \caption{Mapa de amenaza sísmica para SA(0.2), considerando un escenario de ruptura en FSR con $M_w$ 6.8, Dip 60°, caso 3 y profundidad mínima de ruptura de 5 km.}
    \label{fig: Mapa DSHA SA(0.2) FSR 68 CASO 3 DIP 60 PROF MIN 5 KM}
\end{figure}












\begin{figure}
    \centering
    \includegraphics[width=0.7\linewidth]{figuras_tesis/resultados_amenaza/DSHA/SA10 FSR 68 DIP 34 CASO 3 PROF MIN 0 KM.png}
    \caption{Mapa de amenaza sísmica para SA(1.0), considerando un escenario de ruptura en FSR con $M_w$ 6.8, Dip 34°, caso 3 y profundidad mínima de ruptura de 0 km.}
    \label{fig: Mapa DSHA SA(1.0) FSR 68 CASO 3 DIP 34 PROF MIN 0 KM}
\end{figure}

\begin{figure}
    \centering
    \includegraphics[width=0.7\linewidth]{figuras_tesis/resultados_amenaza/DSHA/SA10 FSR 68 DIP 60 CASO 3 PROF MIN 0 KM.png}
    \caption{Mapa de amenaza sísmica para SA(1.0), considerando un escenario de ruptura en FSR con $M_w$ 6.8, Dip 60°, caso 3 y profundidad mínima de ruptura de 0 km.}
    \label{fig: Mapa DSHA SA(1.0) FSR 68 CASO 3 DIP 60 PROF MIN 0 KM}
\end{figure}

\begin{figure}
    \centering
    \includegraphics[width=0.7\linewidth]{figuras_tesis/resultados_amenaza/DSHA/SA10 FSR 68 DIP 34 CASO 3 PROF MIN 5 KM.png}
    \caption{Mapa de amenaza sísmica para SA(1.0), considerando un escenario de ruptura en FSR con $M_w$ 6.8, Dip 34°, caso 3 y profundidad mínima de ruptura de 5 km.}
    \label{fig: Mapa DSHA SA(1.0) FSR 68 CASO 3 DIP 34 PROF MIN 5 KM}
\end{figure}

\begin{figure}
    \centering
    \includegraphics[width=0.7\linewidth]{figuras_tesis/resultados_amenaza/DSHA/SA10 FSR 68 DIP 60 CASO 3 PROF MIN 5 KM.png}
    \caption{Mapa de amenaza sísmica para SA(1.0), considerando un escenario de ruptura en FSR con $M_w$ 6.8, Dip 60°, caso 3 y profundidad mínima de ruptura de 5 km.}
    \label{fig: Mapa DSHA SA(1.0) FSR 68 CASO 3 DIP 60 PROF MIN 5 KM}
\end{figure}








\newpage

\section{Distribuciones de pérdida económica por materialidad y comuna}

En las siguientes secciones, se muestran los resultados de pérdida económica absoluta y relativa al total comunal para cada meterial. Se muestran los resultados para los estudios deterministas y para los estudios basados en catálogos estocásticos de terremotos.

\subsection{Pérdida económica en escenarios sísmicos (DSHA)}

% Definición auxiliar para cabeceras multi-línea
\newcommand{\hdrunit}[2]{\shortstack{#1\\(#2)}}

\subsubsection{Resultados Cabrera et al. (2024)}
%============================ Tabla 1 (abs1) transpuesta ==========================
\begin{table}[H]
	\centering
	\caption{Pérdida económica para escenario de FSR-Poniente $M_w$ 7.5, manteo 34° y profundidad mínima 0 km, para modelo de fragilidad Cabrera et al. (2024).}
	\label{tab: Pérdida escenario FSR-Poniente Cabrera 2024}
	\begin{tabular}{|c|c|c|c|c|}
		\hline
		Comuna       & \hdrunit{Adobe}{CLP $\cdot 10^{6}$} & \hdrunit{RC13}{CLP $\cdot 10^{6}$} & \hdrunit{RM}{CLP $\cdot 10^{6}$} & \hdrunit{Madera}{CLP $\cdot 10^{6}$} \\ \hline
		La Florida   & 0,0    & 0,0     & 30,1     & 4,5    \\ \hline
		La Reina     & 180,5  & 3831,7  & 10686,7  & 660,5  \\ \hline
		Las Condes   & 396,6  & 62031,3 & 48023,8  & 433,1  \\ \hline
		Peñalolén    & 0,0    & 1157,0  & 4047,0   & 65,6   \\ \hline
		Pirque       & 150,3  & 78,1    & 194,2    & 15,1   \\ \hline
		Puente Alto  & 206,8  & 1541,8  & 1584,6   & 65,1   \\ \hline
		Vitacura     & 333,9  & 551,2   & 6593,5   & 0,0    \\ \hline
	\end{tabular}
\end{table}

%============================ Tabla 2 (perc1) transpuesta ==========================
\begin{table}[H]
	\centering
	\caption{Pérdida económica relativa al total por comuna y materialidad para escenario de FSR-Poniente $M_w$ 7.5, manteo 34° y profundidad mínima 0 km, para modelo de fragilidad Cabrera et al. (2024).}
	\label{tab: Pérdida relativa escenario FSR-Poniente Cabrera 2024}
	\begin{tabular}{|c|c|c|c|c|}
		\hline
		Comuna       & Adobe & RC13 & RM  & Madera \\ \hline
		La Florida   & 0\%   & 0\%  & 3\% & 0\%    \\ \hline
		La Reina     & 72\%  & 33\% & 17\%& 6\%    \\ \hline
		Las Condes   & 67\%  & 31\% & 13\%& 4\%    \\ \hline
		Peñalolén    & 0\%   & 31\% & 16\%& 6\%    \\ \hline
		Pirque       & 66\%  & 23\% & 10\%& 2\%    \\ \hline
		Puente Alto  & 70\%  & 32\% & 9\% & 3\%    \\ \hline
		Vitacura     & 68\%  & 25\% & 10\%& 0\%    \\ \hline
	\end{tabular}
\end{table}


% Tabla 1: valores absolutos originales (columnas centradas)
\begin{table}[H]
	\centering
	\caption{Pérdida económica para escenario de FSR-Oriente $M_w$ 7.5, manteo 34° y profundidad mínima 0 km, para modelo de fragilidad Cabrera et al. (2024).}
	\label{tab: Pérdida escenario FSR-Oriente Cabrera 2024}
	\begin{tabular}{|c|c|c|c|c|}
		\hline
		Comuna       & \hdrunit{Adobe}{CLP $\cdot 10^{6}$} & \hdrunit{RC13}{CLP $\cdot 10^{6}$} & \hdrunit{RM}{CLP $\cdot 10^{6}$} & \hdrunit{Madera}{CLP $\cdot 10^{6}$} \\ \hline
		La Florida   & 0,0   & 0,0     & 28,8    & 3,6    \\ \hline
		La Reina     & 175,2 & 3728,7  & 10559,8 & 636,8  \\ \hline
		Las Condes   & 381,3 & 59127,8 & 32473,7 & 198,0  \\ \hline
		Peñalolén    & 0,0   & 1086,7  & 3935,2  & 59,1   \\ \hline
		Pirque       & 151,5 & 88,4    & 205,2   & 15,8   \\ \hline
		Puente Alto  & 211,2 & 1595,6  & 1531,5  & 64,6   \\ \hline
		Vitacura     & 285,6 & 293,4   & 2730,4  & 0,0    \\ \hline
	\end{tabular}
\end{table}

% Tabla 2: porcentajes originales (columnas centradas)
\begin{table}[H]
	\centering
	\caption{Pérdida económica relativa al total por comuna y materialidad para escenario de FSR-Oriente $M_w$ 7.5, manteo 34° y profundidad mínima 0 km, para modelo de fragilidad Cabrera et al. (2024).}
	\label{tab: Pérdida relativa escenario FSR-Oriente Cabrera 2024}
	\begin{tabular}{|c|c|c|c|c|}
		\hline
		Comuna       & Adobe & RC13 & RM  & Madera \\ \hline
		La Florida   & 0\%   & 0\%  & 3\% & 0\%    \\ \hline
		La Reina     & 70\%  & 32\% & 16\%& 6\%    \\ \hline
		Las Condes   & 64\%  & 30\% & 9\% & 2\%    \\ \hline
		Peñalolén    & 0\%   & 29\% & 16\%& 5\%    \\ \hline
		Pirque       & 67\%  & 26\% & 10\%& 2\%    \\ \hline
		Puente Alto  & 72\%  & 33\% & 9\% & 3\%    \\ \hline
		Vitacura     & 59\%  & 13\% & 4\% & 0\%    \\ \hline
	\end{tabular}
\end{table}

% Tabla 3: valores absolutos actualizados (columnas centradas)
\begin{table}[H]
	\centering
	\caption{Pérdida económica para escenario de FSR-Multitraza $M_w$ 7.5, manteo 34° y profundidad mínima 0 km, para modelo de fragilidad Cabrera et al. (2024).}
	\label{tab: Pérdida escenario FSR-Multitraza Cabrera 2024}
	\begin{tabular}{|c|c|c|c|c|}
		\hline
		Comuna       & \hdrunit{Adobe}{CLP $\cdot 10^{6}$} & \hdrunit{RC13}{CLP $\cdot 10^{6}$} & \hdrunit{RM}{CLP $\cdot 10^{6}$} & \hdrunit{Madera}{CLP $\cdot 10^{6}$} \\ \hline
		La Florida   & 0,0   & 0,0     & 28,8    & 3,5    \\ \hline
		La Reina     & 175,2 & 3956,0  & 10666,2 & 644,0  \\ \hline
		Las Condes   & 410,0 & 67238,8 & 52117,0 & 437,8  \\ \hline
		Peñalolén    & 0,0   & 1086,7  & 3934,2  & 59,1   \\ \hline
		Pirque       & 151,5 & 88,3    & 204,9   & 15,8   \\ \hline
		Puente Alto  & 211,2 & 1594,7  & 1531,2  & 64,6   \\ \hline
		Vitacura     & 356,6 & 558,0   & 7582,9  & 0,0    \\ \hline
	\end{tabular}
\end{table}

% Tabla 4: porcentajes actualizados (columnas centradas)
\begin{table}[H]
	\centering
	\caption{Pérdida económica relativa al total por comuna y materialidad para escenario de FSR-Multitraza $M_w$ 7.5, manteo 34° y profundidad mínima 0 km, para modelo de fragilidad Cabrera et al. (2024).}
	\label{tab: Pérdida relativa escenario FSR-Multitraza Cabrera 2024}
	\begin{tabular}{|c|c|c|c|c|}
		\hline
		Comuna       & Adobe & RC13 & RM  & Madera \\ \hline
		La Florida   & 0\%   & 0\%  & 3\% & 0\%    \\ \hline
		La Reina     & 70\%  & 34\% & 17\%& 6\%    \\ \hline
		Las Condes   & 69\%  & 34\% & 14\%& 5\%    \\ \hline
		Peñalolén    & 0\%   & 29\% & 16\%& 5\%    \\ \hline
		Pirque       & 67\%  & 26\% & 10\%& 2\%    \\ \hline
		Puente Alto  & 72\%  & 33\% & 9\% & 3\%    \\ \hline
		Vitacura     & 73\%  & 25\% & 12\%& 0\%    \\ \hline
	\end{tabular}
\end{table}


%============================ Tabla 3 (abs2) transpuesta ==========================
\begin{table}[H]
	\centering
	\caption{Pérdida económica para escenario de Interplaca $M_w$ 9.3, para modelo de fragilidad Cabrera et al. (2024).}
	\label{tab: Pérdida escenario Interplaca Cabrera 2024}
	\begin{tabular}{|c|c|c|c|c|}
		\hline
		Comuna       & \hdrunit{Adobe}{CLP $\cdot 10^{6}$} & \hdrunit{RC13}{CLP $\cdot 10^{6}$} & \hdrunit{RM}{CLP $\cdot 10^{6}$} & \hdrunit{Madera}{CLP $\cdot 10^{6}$} \\ \hline
		La Florida   & 0,0    & 0,0    & 3,6    & 0,2    \\ \hline
		La Reina     & 86,5   & 315,7  & 488,3  & 5,7    \\ \hline
		Las Condes   & 195,1  & 5133,9 & 2728,2 & 5,5    \\ \hline
		Peñalolén    & 0,0    & 100,8  & 188,8  & 0,6    \\ \hline
		Pirque       & 75,3   & 10,1   & 16,7   & 0,5    \\ \hline
		Puente Alto  & 101,6  & 133,9  & 116,6  & 1,6    \\ \hline
		Vitacura     & 171,5  & 71,5   & 575,7  & 0,0    \\ \hline
	\end{tabular}
\end{table}

%============================ Tabla 4 (perc2) transpuesta ==========================
\begin{table}[H]
	\centering
	\caption{Pérdida económica relativa al total por material y por comuna para escenario de Interplaca $M_w$ 9.3, para modelo de fragilidad Cabrera et al. (2024).}
	\label{tab: Pérdida relativa escenario Interplaca Cabrera 2024}
	\begin{tabular}{|c|c|c|c|c|}
		\hline
		Comuna       & Adobe & RC13 & RM  & Madera \\ \hline
		La Florida   & 0\%   & 0\%  & 0\% & 0\%    \\ \hline
		La Reina     & 35\%  & 3\%  & 1\% & 0\%    \\ \hline
		Las Condes   & 33\%  & 3\%  & 1\% & 0\%    \\ \hline
		Peñalolén    & 0\%   & 3\%  & 1\% & 0\%    \\ \hline
		Pirque       & 33\%  & 3\%  & 1\% & 0\%    \\ \hline
		Puente Alto  & 35\%  & 3\%  & 1\% & 0\%    \\ \hline
		Vitacura     & 35\%  & 3\%  & 1\% & 0\%    \\ \hline
	\end{tabular}
\end{table}

%============================ Tabla 5 (abs3) transpuesta ==========================
\begin{table}[H]
	\centering
	\caption{Pérdida económica para escenario de Intraplaca $M_w$ 8.5, para modelo de fragilidad Cabrera et al. (2024).}
	\label{tab: Pérdida escenario Intraplaca Cabrera 2024}
	\begin{tabular}{|c|c|c|c|c|}
		\hline
		Comuna       & \hdrunit{Adobe}{CLP $\cdot 10^{6}$} & \hdrunit{RC13}{CLP $\cdot 10^{6}$} & \hdrunit{RM}{CLP $\cdot 10^{6}$} & \hdrunit{Madera}{CLP $\cdot 10^{6}$} \\ \hline
		La Florida   & 0,0    & 0,0     & 13,5   & 2,2    \\ \hline
		La Reina     & 78,7   & 548,4   & 1024,6 & 24,6   \\ \hline
		Las Condes   & 179,0  & 9500,0  & 6529,2 & 26,8   \\ \hline
		Peñalolén    & 0,0    & 174,2   & 394,6  & 2,5    \\ \hline
		Pirque       & 68,7   & 22,9    & 49,3   & 3,4    \\ \hline
		Puente Alto  & 92,0   & 234,8   & 314,6  & 9,1    \\ \hline
		Vitacura     & 155,5  & 139,9   & 1515,0 & 0,0    \\ \hline
	\end{tabular}
\end{table}

%============================ Tabla 6 (perc3) transpuesta ==========================
\begin{table}[H]
	\centering
	\caption{Pérdida económica relativa al total por material y por comuna para escenario de Intraplaca $M_w$ 8.5, para modelo de fragilidad Cabrera et al. (2024).}
	\label{tab: Pérdida relativa escenario Intraplaca Cabrera 2024}
	\begin{tabular}{|c|c|c|c|c|}
		\hline
		Comuna       & Adobe & RC13 & RM  & Madera \\ \hline
		La Florida   & 0\%   & 0\%  & 1\% & 0\%    \\ \hline
		La Reina     & 31\%  & 5\%  & 2\% & 0\%    \\ \hline
		Las Condes   & 30\%  & 5\%  & 2\% & 0\%    \\ \hline
		Peñalolén    & 0\%   & 5\%  & 2\% & 0\%    \\ \hline
		Pirque       & 30\%  & 7\%  & 2\% & 0\%    \\ \hline
		Puente Alto  & 31\%  & 5\%  & 2\% & 0\%    \\ \hline
		Vitacura     & 32\%  & 6\%  & 2\% & 0\%    \\ \hline
	\end{tabular}
\end{table}

\subsubsection{Resultados Favier et al. (2017)}

%============================ Tabla 1 (abs1) transpuesta ==========================
\begin{table}[H]
	\centering
	\caption{Pérdida económica para escenario de FSR-Poniente $M_w$ 7.5, manteo 34° y profundidad mínima 0 km, para modelo de fragilidad Favier et al. (2017).}
	\label{tab: Pérdida escenario FSR-Poniente Favier 2017}
	\begin{tabular}{|c|c|c|c|}
		\hline
		Comuna       & \hdrunit{RC1024}{CLP $\cdot 10^{6}$} & \hdrunit{RC13}{CLP $\cdot 10^{6}$} & \hdrunit{RC49}{CLP $\cdot 10^{6}$} \\ \hline
		La Reina     & 0,0    & 10031,5  & 0,0    \\ \hline
		Las Condes   & 5690,6 & 170836,5 & 89971,3 \\ \hline
		Peñalolén    & 0,0    & 3169,5   & 14828,9 \\ \hline
		Pirque       & 0,0    & 272,7    & 0,0     \\ \hline
		Puente Alto  & 0,0    & 4115,6   & 0,0     \\ \hline
		Vitacura     & 108322,4 & 1845,0 & 6669,3  \\ \hline
	\end{tabular}
\end{table}

%============================ Tabla 2 (perc1) transpuesta ==========================
\begin{table}[H]
	\centering
	\caption{Pérdida económica relativa al total por comuna y materialidad para escenario de FSR-Poniente $M_w$ 7.5, manteo 34° y profundidad mínima 0 km, para modelo de fragilidad Favier et al. (2017).}
	\label{tab: Pérdida relativa escenario FSR-Poniente Favier 2017}
	\begin{tabular}{|c|c|c|c|}
		\hline
		Comuna       & RC1024 & RC13 & RC49 \\ \hline
		La Reina     & 0\%  & 86\% & 0\%  \\ \hline
		Las Condes   & 99\% & 85\% & 93\% \\ \hline
		Peñalolén    & 0\%  & 85\% & 99\% \\ \hline
		Pirque       & 0\%  & 80\% & 0\%  \\ \hline
		Puente Alto  & 0\%  & 86\% & 0\%  \\ \hline
		Vitacura     & 99\% & 83\% & 99\% \\ \hline
	\end{tabular}
\end{table}













% Tabla 1 transpuesta: valores absolutos originales
\begin{table}[H]
	\centering
	\caption{Pérdida económica para escenario de FSR-Oriente $M_w$ 7.5, manteo 34° y profundidad mínima 0 km, para modelo de fragilidad Favier et al. (2017).}
	\label{tab: Pérdida escenario FSR-Oriente Favier 2017}	
	\begin{tabular}{|c|c|c|c|}
		\hline
		Comuna       & \hdrunit{RC1024}{CLP $\cdot 10^{6}$} & \hdrunit{RC13}{CLP $\cdot 10^{6}$} & \hdrunit{RC49}{CLP $\cdot 10^{6}$} \\ \hline
		La Reina      & 0,0       & 10148,0   & 0,0       \\ \hline
		Las Condes    & 5647,3    & 169053,1  & 90291,5   \\ \hline
		Peñalolén     & 0,0       & 3133,0    & 14870,3   \\ \hline
		Pirque        & 0,0       & 290,8     & 0,0       \\ \hline
		Puente Alto   & 0,0       & 4267,5    & 0,0       \\ \hline
		Vitacura      & 106707,1  & 1492,3    & 6563,1    \\ \hline
	\end{tabular}
\end{table}

% Tabla 2 transpuesta: porcentajes originales
\begin{table}[H]
	\centering
	\caption{Pérdida económica relativa al total por comuna y materialidad para escenario de FSR-Oriente $M_w$ 7.5, manteo 34° y profundidad mínima 0 km, para modelo de fragilidad Favier et al. (2017).}
	\label{tab: Pérdida relativa escenario FSR-Oriente Favier 2017}
	\begin{tabular}{|c|c|c|c|}
		\hline
		Comuna       & RC1024 & RC13 & RC49 \\ \hline
		La Reina      & 0\%    & 87\%  & 0\%   \\ \hline
		Las Condes    & 99\%   & 85\%  & 93\%  \\ \hline
		Peñalolén     & 0\%    & 84\%  & 100\% \\ \hline
		Pirque        & 0\%    & 85\%  & 0\%   \\ \hline
		Puente Alto   & 0\%    & 89\%  & 0\%   \\ \hline
		Vitacura      & 98\%   & 67\%  & 98\%  \\ \hline
	\end{tabular}
\end{table}

% Tabla 3 transpuesta: valores absolutos actualizados
\begin{table}[H]
	\centering
	\caption{Pérdida económica para escenario de FSR-Multitraza $M_w$ 7.5, manteo 34° y profundidad mínima 0 km, para modelo de fragilidad Favier et al. (2017).}
	\label{tab: Pérdida escenario FSR-Multitraza Favier 2017}
	\begin{tabular}{|c|c|c|c|}
		\hline
		Comuna       & \hdrunit{RC1024}{CLP $\cdot 10^{6}$} & \hdrunit{RC13}{CLP $\cdot 10^{6}$} & \hdrunit{RC49}{CLP $\cdot 10^{6}$} \\ \hline
		La Reina      & 0,0       & 10425,4    & 0,0       \\ \hline
		Las Condes    & 5699,6    & 178107,9   & 91609,3   \\ \hline
		Peñalolén     & 0,0       & 3133,0     & 14870,3   \\ \hline
		Pirque        & 0,0       & 290,7      & 0,0       \\ \hline
		Puente Alto   & 0,0       & 4266,8     & 0,0       \\ \hline
		Vitacura      & 108439,0  & 1856,4     & 6694,3    \\ \hline
	\end{tabular}
\end{table}

% Tabla 4 transpuesta: porcentajes actualizados
\begin{table}[H]
	\centering
	\caption{Pérdida económica relativa al total por comuna y materialidad para escenario de FSR-Multitraza $M_w$ 7.5, manteo 34° y profundidad mínima 0 km, para modelo de fragilidad Favier et al. (2017).}
	\label{tab: Pérdida relativa escenario FSR-Multitraza Favier 2017}
	\begin{tabular}{|c|c|c|c|}
		\hline
		Comuna       & RC1024 & RC13 & RC49 \\ \hline
		La Reina      & 0\%    & 89\%  & 0\%   \\ \hline
		Las Condes    & 100\%  & 89\%  & 95\%  \\ \hline
		Peñalolén     & 0\%    & 84\%  & 100\% \\ \hline
		Pirque        & 0\%    & 85\%  & 0\%   \\ \hline
		Puente Alto   & 0\%    & 89\%  & 0\%   \\ \hline
		Vitacura      & 99\%   & 83\%  & 100\% \\ \hline
	\end{tabular}
\end{table}











%============================ Tabla 3 (abs2) transpuesta ==========================
\begin{table}[H]
	\centering
	\caption{Pérdida económica para escenario de Interplaca $M_w$ 9.3, para modelo de fragilidad Favier et al. (2017).}
	\label{tab: Pérdida escenario Interplaca Favier 2017}
	\begin{tabular}{|c|c|c|c|}
		\hline
		Comuna       & \hdrunit{RC1024}{CLP $\cdot 10^{6}$} & \hdrunit{RC13}{CLP $\cdot 10^{6}$} & \hdrunit{RC49}{CLP $\cdot 10^{6}$} \\ \hline
		La Reina     & 0,0    & 2294,5   & 0,0     \\ \hline
		Las Condes   & 4087,3 & 37059,2  & 55691,5 \\ \hline
		Peñalolén    & 0,0    & 732,1    & 10278,5 \\ \hline
		Pirque       & 0,0    & 71,6     & 0,0     \\ \hline
		Puente Alto  & 0,0    & 971,3    & 0,0     \\ \hline
		Vitacura     & 78047,4 & 506,8   & 4518,4  \\ \hline
	\end{tabular}
\end{table}

%============================ Tabla 4 (perc2) transpuesta ==========================
\begin{table}[H]
	\centering
	\caption{Pérdida económica relativa al total por material y por comuna para escenario de Interplaca $M_w$ 9.3, para modelo de fragilidad Favier et al. (2017).}
	\label{tab: Pérdida relativa escenario Interplaca Favier 2017}
	\begin{tabular}{|c|c|c|c|}
		\hline
		Comuna       & RC1024 & RC13 & RC49 \\ \hline
		La Reina     & 0\%  & 20\% & 0\%  \\ \hline
		Las Condes   & 71\% & 19\% & 58\% \\ \hline
		Peñalolén    & 0\%  & 20\% & 69\% \\ \hline
		Pirque       & 0\%  & 21\% & 0\%  \\ \hline
		Puente Alto  & 0\%  & 20\% & 0\%  \\ \hline
		Vitacura     & 72\% & 23\% & 67\% \\ \hline
	\end{tabular}
\end{table}

%============================ Tabla 5 (abs3) transpuesta ==========================
\begin{table}[H]
	\centering
	\caption{Pérdida económica para escenario de Intraplaca $M_w$ 8.5, para modelo de fragilidad Favier et al. (2017).}
	\label{tab: Pérdida escenario Intraplaca Favier 2017}
	\begin{tabular}{|c|c|c|c|}
		\hline
		Comuna       & \hdrunit{RC1024}{CLP $\cdot 10^{6}$} & \hdrunit{RC13}{CLP $\cdot 10^{6}$} & \hdrunit{RC49}{CLP $\cdot 10^{6}$} \\ \hline
		La Reina     & 0,0    & 2391,3   & 0,0     \\ \hline
		Las Condes   & 2763,4 & 40343,8  & 36268,9 \\ \hline
		Peñalolén    & 0,0    & 762,2    & 6968,1  \\ \hline
		Pirque       & 0,0    & 71,1     & 0,0     \\ \hline
		Puente Alto  & 0,0    & 989,7    & 0,0     \\ \hline
		Vitacura     & 53027,4 & 474,4   & 2968,5  \\ \hline
	\end{tabular}
\end{table}

%============================ Tabla 6 (perc3) transpuesta ==========================
\begin{table}[H]
	\centering
	\caption{Pérdida económica relativa al total por material y por comuna para escenario de Intraplaca $M_w$ 8.5, para modelo de fragilidad Favier et al. (2017).}
	\label{tab: Pérdida relativa escenario Intraplaca Favier 2017}
	\begin{tabular}{|c|c|c|c|}
		\hline
		Comuna       & RC1024 & RC13 & RC49 \\ \hline
		La Reina     & 0\%  & 20\% & 0\%  \\ \hline
		Las Condes   & 48\% & 20\% & 38\% \\ \hline
		Peñalolén    & 0\%  & 20\% & 47\% \\ \hline
		Pirque       & 0\%  & 21\% & 0\%  \\ \hline
		Puente Alto  & 0\%  & 21\% & 0\%  \\ \hline
		Vitacura     & 49\% & 21\% & 44\% \\ \hline
	\end{tabular}
\end{table}

\subsubsection{Resultados HAZUS (2024)}

%% =======================================================================
%% Conjunto 1
%% =======================================================================

\begin{table}[H]
	\centering
	\caption{Pérdida económica para escenario de FSR-Poniente $M_w$ 7.5, manteo 34° y profundidad mínima 0 km, para modelo de fragilidad HAZUS (2024).}
	\label{tab: Pérdida escenario FSR-Poniente HAZUS 2024}
	\begin{tabular}{|c|c|c|c|c|c|}
		\hline
		Comuna       & \hdrunit{RC13}{CLP $\cdot 10^{6}$} & \hdrunit{RC47}{CLP $\cdot 10^{6}$} & \hdrunit{RC8}{CLP $\cdot 10^{6}$}  & \hdrunit{RM}{CLP $\cdot 10^{6}$}   & \hdrunit{Madera}{CLP $\cdot 10^{6}$} \\ \hline
		La Florida   & 0,0    & 0,0     & 0,0      & 353,4    & 285,6    \\ \hline
		La Reina     & 6647,4 & 0,0     & 0,0      & 37552,9  & 5165,6   \\ \hline
		Las Condes   & 111721,6 & 48609,5 & 3378,8  & 204456,2 & 4307,2   \\ \hline
		Peñalolén    & 2079,0 & 8826,2  & 0,0      & 14434,9  & 517,3    \\ \hline
		Pirque       & 174,5  & 0,0     & 0,0      & 1065,7   & 343,9    \\ \hline
		Puente Alto  & 2715,8 & 0,0     & 0,0      & 8444,3   & 1079,2   \\ \hline
		Vitacura     & 1165,0 & 3719,9  & 64340,4  & 34445,8  & 0,0      \\ \hline
	\end{tabular}
\end{table}

\begin{table}[H]
	\centering
	\caption{Pérdida económica relativa al total por comuna y materialidad para escenario de FSR-Poniente $M_w$ 7.5, manteo 34° y profundidad mínima 0 km, para modelo de fragilidad HAZUS (2024).}
	\label{tab: Pérdida relativa escenario FSR-Poniente HAZUS 2024}
	\begin{tabular}{|c|c|c|c|c|c|}
		\hline
		Comuna       & RC13 & RC47 & RC8  & RM  & Madera \\ \hline
		La Florida   & 0\%  & 0\%  & 0\%  & 39\% & 29\%  \\ \hline
		La Reina     & 57\% & 0\%  & 0\%  & 58\% & 47\%  \\ \hline
		Las Condes   & 56\% & 50\% & 59\% & 54\% & 45\%  \\ \hline
		Peñalolén    & 56\% & 59\% & 0\%  & 58\% & 45\%  \\ \hline
		Pirque       & 51\% & 0\%  & 0\%  & 54\% & 40\%  \\ \hline
		Puente Alto  & 57\% & 0\%  & 0\%  & 49\% & 43\%  \\ \hline
		Vitacura     & 52\% & 55\% & 59\% & 54\% & 0\%   \\ \hline
	\end{tabular}
\end{table}












% Tabla 1 transpuesta: valores absolutos originales
\begin{table}[H]
	\centering
	\caption{Pérdida económica para escenario de FSR-Oriente $M_w$ 7.5, manteo 34° y profundidad mínima 0 km, para modelo de fragilidad HAZUS (2024).}
	\label{tab: Pérdida escenario FSR-Oriente HAZUS 2024}
	\begin{tabular}{|c|c|c|c|c|c|}
		\hline
		Comuna       & \hdrunit{RC13}{CLP $\cdot 10^{6}$} & \hdrunit{RC47}{CLP $\cdot 10^{6}$} & \hdrunit{RC8}{CLP $\cdot 10^{6}$}  & \hdrunit{RM}{CLP $\cdot 10^{6}$}   & \hdrunit{Madera}{CLP $\cdot 10^{6}$} \\ \hline
		La Florida    & 0,0        & 0,0        & 0,0        & 356,1      & 276,8      \\ \hline
		La Reina      & 6611,8     & 0,0        & 0,0        & 37876,5    & 5167,4     \\ \hline
		Las Condes    & 109798,1   & 47196,8    & 3057,6     & 188347,1   & 3812,3     \\ \hline
		Peñalolén     & 2032,3     & 8986,2     & 0,0        & 14444,6    & 509,4      \\ \hline
		Pirque        & 184,3      & 0,0        & 0,0        & 1096,5     & 356,4      \\ \hline
		Puente Alto   & 2779,0     & 0,0        & 0,0        & 8386,9     & 1098,4     \\ \hline
		Vitacura      & 951,0      & 3235,9     & 56346,4    & 28396,6    & 0,0        \\ \hline
	\end{tabular}
\end{table}

% Tabla 2 transpuesta: porcentajes originales
\begin{table}[H]
	\centering
	\caption{Pérdida económica relativa al total por comuna y materialidad para escenario de FSR-Oriente $M_w$ 7.5, manteo 34° y profundidad mínima 0 km, para modelo de fragilidad HAZUS (2024).}
	\label{tab: Pérdida relativa escenario FSR-Oriente HAZUS 2024}
	\begin{tabular}{|c|c|c|c|c|c|}
		\hline
		Comuna       & RC13 & RC47 & RC8  & RM  & Madera \\ \hline
		La Florida    & 0\%  & 0\%  & 0\%  & 39\% & 28\%   \\ \hline
		La Reina      & 57\% & 0\%  & 0\%  & 59\% & 47\%   \\ \hline
		Las Condes    & 55\% & 49\% & 53\% & 49\% & 39\%   \\ \hline
		Peñalolén     & 54\% & 60\% & 0\%  & 58\% & 44\%   \\ \hline
		Pirque        & 54\% & 0\%  & 0\%  & 55\% & 41\%   \\ \hline
		Puente Alto   & 58\% & 0\%  & 0\%  & 49\% & 44\%   \\ \hline
		Vitacura      & 43\% & 48\% & 52\% & 44\% & 0\%    \\ \hline
	\end{tabular}
\end{table}

% Tabla 3 transpuesta: valores absolutos actualizados
\begin{table}[H]
	\centering
	\caption{Pérdida económica para escenario de FSR-Multitraza $M_w$ 7.5, manteo 34° y profundidad mínima 0 km, para modelo de fragilidad HAZUS (2024).}
	\label{tab: Pérdida escenario FSR-Multitraza HAZUS 2024}
	\begin{tabular}{|c|c|c|c|c|c|}
		\hline
		Comuna       & \hdrunit{RC13}{CLP $\cdot 10^{6}$} & \hdrunit{RC47}{CLP $\cdot 10^{6}$} & \hdrunit{RC8}{CLP $\cdot 10^{6}$}  & \hdrunit{RM}{CLP $\cdot 10^{6}$}   & \hdrunit{Madera}{CLP $\cdot 10^{6}$} \\ \hline
		La Florida    & 0,0        & 0,0        & 0,0        & 356,0      & 276,8      \\ \hline
		La Reina      & 6802,7     & 0,0        & 0,0        & 38059,7    & 5216,0     \\ \hline
		Las Condes    & 117121,2   & 51063,5    & 3557,3     & 214935,1   & 4516,0     \\ \hline
		Peñalolén     & 2032,3     & 8984,2     & 0,0        & 14443,1    & 509,4      \\ \hline
		Pirque        & 184,2      & 0,0        & 0,0        & 1096,2     & 356,2      \\ \hline
		Puente Alto   & 2778,3     & 0,0        & 0,0        & 8386,5     & 1098,1     \\ \hline
		Vitacura      & 1172,1     & 3954,5     & 66544,9    & 36204,3    & 0,0        \\ \hline
	\end{tabular}
\end{table}

% Tabla 4 transpuesta: porcentajes actualizados
\begin{table}[H]
	\centering
	\caption{Pérdida económica relativa al total por comuna y materialidad para escenario de FSR-Multitraza $M_w$ 7.5, manteo 34° y profundidad mínima 0 km, para modelo de fragilidad HAZUS (2024).}
	\label{tab: Pérdida relativa escenario FSR-Multitraza HAZUS 2024}
	\begin{tabular}{|c|c|c|c|c|c|}
		\hline
		Comuna       & RC13 & RC47 & RC8  & RM  & Madera \\ \hline
		La Florida    & 0\%  & 0\%  & 0\%  & 39\% & 28\%   \\ \hline
		La Reina      & 58\% & 0\%  & 0\%  & 59\% & 47\%   \\ \hline
		Las Condes    & 59\% & 53\% & 62\% & 56\% & 47\%   \\ \hline
		Peñalolén     & 54\% & 60\% & 0\%  & 58\% & 44\%   \\ \hline
		Pirque        & 54\% & 0\%  & 0\%  & 55\% & 41\%   \\ \hline
		Puente Alto   & 58\% & 0\%  & 0\%  & 49\% & 44\%   \\ \hline
		Vitacura      & 53\% & 59\% & 61\% & 57\% & 0\%    \\ \hline
	\end{tabular}
\end{table}










%% =======================================================================
%% Conjunto 2
%% =======================================================================

\begin{table}[H]
	\centering
	\caption{Pérdida económica para escenario de Interplaca $M_w$ 9.3, para modelo de fragilidad HAZUS (2024).}
	\label{tab: Pérdida escenario Interplaca HAZUS 2024}
	\begin{tabular}{|c|c|c|c|c|c|}
		\hline
		Comuna       & \hdrunit{RC13}{CLP $\cdot 10^{6}$} & \hdrunit{RC47}{CLP $\cdot 10^{6}$} & \hdrunit{RC8}{CLP $\cdot 10^{6}$}  & \hdrunit{RM}{CLP $\cdot 10^{6}$}   & \hdrunit{Madera}{CLP $\cdot 10^{6}$} \\ \hline
		La Florida   & 0,0    & 0,0     & 0,0      & 171,1    & 128,8    \\ \hline
		La Reina     & 2790,5 & 0,0     & 0,0      & 16310,6  & 1962,6   \\ \hline
		Las Condes   & 46269,8 & 26176,0 & 2110,5  & 92091,2  & 1715,3   \\ \hline
		Peñalolén    & 893,0  & 4531,2  & 0,0      & 6321,9   & 203,1    \\ \hline
		Pirque       & 78,9   & 0,0     & 0,0      & 488,5    & 143,7    \\ \hline
		Puente Alto  & 1153,8 & 0,0     & 0,0      & 3986,7   & 439,7    \\ \hline
		Vitacura     & 543,8  & 2069,2  & 40268,9  & 16442,4  & 0,0      \\ \hline
	\end{tabular}
\end{table}

\begin{table}[H]
	\centering
	\caption{Pérdida económica relativa al total por material y por comuna para escenario de Interplaca $M_w$ 9.3, para modelo de fragilidad HAZUS (2024).}
	\label{tab: Pérdida relativa escenario Interplaca HAZUS 2024}
	\begin{tabular}{|c|c|c|c|c|c|}
		\hline
		Comuna       & RC13 & RC47 & RC8  & RM  & Madera \\ \hline
		La Florida   & 0\%  & 0\%  & 0\%  & 19\% & 13\%  \\ \hline
		La Reina     & 24\% & 0\%  & 0\%  & 25\% & 18\%  \\ \hline
		Las Condes   & 23\% & 27\% & 37\% & 24\% & 18\%  \\ \hline
		Peñalolén    & 24\% & 30\% & 0\%  & 25\% & 18\%  \\ \hline
		Pirque       & 23\% & 0\%  & 0\%  & 25\% & 17\%  \\ \hline
		Puente Alto  & 24\% & 0\%  & 0\%  & 23\% & 18\%  \\ \hline
		Vitacura     & 24\% & 31\% & 37\% & 26\% & 0\%   \\ \hline
	\end{tabular}
\end{table}


%% =======================================================================
%% Conjunto 3
%% =======================================================================

\begin{table}[H]
	\centering
	\caption{Pérdida económica para escenario de Intraplaca $M_w$ 8.5, para modelo de fragilidad HAZUS (2024).}
	\label{tab: Pérdida escenario Intraplaca HAZUS 2024}
	\begin{tabular}{|c|c|c|c|c|c|}
		\hline
		Comuna       & \hdrunit{RC13}{CLP $\cdot 10^{6}$} & \hdrunit{RC47}{CLP $\cdot 10^{6}$} & \hdrunit{RC8}{CLP $\cdot 10^{6}$}  & \hdrunit{RM}{CLP $\cdot 10^{6}$}   & \hdrunit{Madera}{CLP $\cdot 10^{6}$} \\ \hline
		La Florida   & 0,0    & 0,0     & 0,0      & 165,5    & 131,5    \\ \hline
		La Reina     & 2565,3 & 0,0     & 0,0      & 14872,6  & 1838,3   \\ \hline
		Las Condes   & 42793,9 & 23978,8 & 1928,8  & 84734,0  & 1615,7   \\ \hline
		Peñalolén    & 818,3  & 4113,7  & 0,0      & 5751,8   & 189,6    \\ \hline
		Pirque       & 74,0   & 0,0     & 0,0      & 450,2    & 139,2    \\ \hline
		Puente Alto  & 1058,6 & 0,0     & 0,0      & 3694,7   & 418,2    \\ \hline
		Vitacura     & 503,1  & 1878,0  & 36845,2  & 15059,2  & 0,0      \\ \hline
	\end{tabular}
\end{table}

\begin{table}[H]
	\centering
	\caption{Pérdida económica relativa al total por material y por comuna para escenario de Intraplaca $M_w$ 8.5, para modelo de fragilidad HAZUS (2024).}
	\label{tab: Pérdida relativa escenario Intraplaca HAZUS 2024}
	\begin{tabular}{|c|c|c|c|c|c|}
		\hline
		Comuna       & RC13 & RC47 & RC8  & RM  & Madera \\ \hline
		La Florida   & 0\%  & 0\%  & 0\%  & 18\% & 13\%  \\ \hline
		La Reina     & 22\% & 0\%  & 0\%  & 23\% & 17\%  \\ \hline
		Las Condes   & 21\% & 25\% & 34\% & 22\% & 17\%  \\ \hline
		Peñalolén    & 22\% & 28\% & 0\%  & 23\% & 17\%  \\ \hline
		Pirque       & 22\% & 0\%  & 0\%  & 23\% & 16\%  \\ \hline
		Puente Alto  & 22\% & 0\%  & 0\%  & 21\% & 17\%  \\ \hline
		Vitacura     & 23\% & 28\% & 34\% & 24\% & 0\%   \\ \hline
	\end{tabular}
\end{table}


%% =======================================================================
%% Conjunto 4
%% =======================================================================

\subsection{Pérdida económica basada en catálogo estocástico de terremotos}

\subsubsection{Resultados Cabrera et al. (2024)}

%============================ Tabla 7 (abs4) transpuesta ==========================
\begin{table}[H]
	\centering
	\caption{Pérdida anual esperada (AAL) para fuente FSR, usando modelo de fragilidad Cabrera et al. (2024).}
	\label{tab: AAL FSR Cabrera 2024}
	\begin{tabular}{|c|c|c|c|c|}
		\hline
		Comuna       & \hdrunit{Adobe}{CLP $\cdot 10^{6}$} & \hdrunit{RC13}{CLP $\cdot 10^{6}$} & \hdrunit{RM}{CLP $\cdot 10^{6}$} & \hdrunit{Madera}{CLP $\cdot 10^{6}$} \\ \hline
		La Florida   & 0,000  & 0,000   & 0,009   & 0,004  \\ \hline
		La Reina     & 0,021  & 0,422   & 1,222   & 0,084  \\ \hline
		Las Condes   & 0,046  & 6,824   & 6,098   & 0,066  \\ \hline
		Peñalolén    & 0,000  & 0,128   & 0,459   & 0,008  \\ \hline
		Pirque       & 0,018  & 0,011   & 0,037   & 0,006  \\ \hline
		Puente Alto  & 0,024  & 0,174   & 0,255   & 0,019  \\ \hline
		Vitacura     & 0,039  & 0,073   & 1,102   & 0,000  \\ \hline
	\end{tabular}
\end{table}

%============================ Tabla 8 (perc4) transpuesta ==========================
\begin{table}[H]
	\centering
	\caption{Pérdida anual esperada relativa al total por comuna y por materialidad (AALR) para fuente FSR, usando modelo de fragilidad Cabrera et al. (2024).}
	\label{tab: AALR FSR Cabrera 2024}
	\begin{tabular}{|c|c|c|c|c|}
		\hline
		Comuna       & Adobe    & RC13     & RM       & Madera   \\ \hline
		La Florida   & 0,000\%  & 0,000\%  & 0,001\%  & 0,000\%  \\ \hline
		La Reina     & 0,008\%  & 0,004\%  & 0,002\%  & 0,001\%  \\ \hline
		Las Condes   & 0,008\%  & 0,004\%  & 0,002\%  & 0,001\%  \\ \hline
		Peñalolén    & 0,000\%  & 0,003\%  & 0,002\%  & 0,001\%  \\ \hline
		Pirque       & 0,008\%  & 0,003\%  & 0,002\%  & 0,001\%  \\ \hline
		Puente Alto  & 0,008\%  & 0,004\%  & 0,002\%  & 0,001\%  \\ \hline
		Vitacura     & 0,008\%  & 0,003\%  & 0,002\%  & 0,000\%  \\ \hline
	\end{tabular}
\end{table}

%============================ Tabla 9 (abs5) transpuesta ==========================
\begin{table}[H]
	\centering
	\caption{Pérdida anual esperada (AAL) para fuente Interplaca, usando modelo de fragilidad Cabrera et al. (2024).}
	\label{tab: AAL Interplaca Cabrera 2024}
	\begin{tabular}{|c|c|c|c|c|}
		\hline
		Comuna       & \hdrunit{Adobe}{CLP $\cdot 10^{6}$} & \hdrunit{RC13}{CLP $\cdot 10^{6}$} & \hdrunit{RM}{CLP $\cdot 10^{6}$} & \hdrunit{Madera}{CLP $\cdot 10^{6}$} \\ \hline
		La Florida   & 0,000  & 0,000   & 0,137   & 0,045   \\ \hline
		La Reina     & 0,890  & 6,640   & 16,316  & 0,872   \\ \hline
		Las Condes   & 2,053  & 106,722 & 89,802  & 0,774   \\ \hline
		Peñalolén    & 0,000  & 2,091   & 6,116   & 0,098   \\ \hline
		Pirque       & 0,707  & 0,170   & 0,453   & 0,059   \\ \hline
		Puente Alto  & 0,963  & 2,640   & 3,670   & 0,184   \\ \hline
		Vitacura     & 1,730  & 1,317   & 17,263  & 0,000   \\ \hline
	\end{tabular}
\end{table}

%============================ Tabla 10 (perc5) transpuesta ==========================
\begin{table}[H]
	\centering
	\caption{Pérdida anual esperada relativa al total por comuna y por materialidad (AALR) para fuente Interplaca, usando modelo de fragilidad Cabrera et al. (2024).}
	\label{tab: AALR Interplaca Cabrera 2024}
	\begin{tabular}{|c|c|c|c|c|}
		\hline
		Comuna       & Adobe    & RC13     & RM       & Madera   \\ \hline
		La Florida   & 0,000\%  & 0,000\%  & 0,015\%  & 0,005\%  \\ \hline
		La Reina     & 0,356\%  & 0,057\%  & 0,025\%  & 0,008\%  \\ \hline
		Las Condes   & 0,332\%  & 0,055\%  & 0,025\%  & 0,008\%  \\ \hline
		Peñalolén    & 0,000\%  & 0,055\%  & 0,024\%  & 0,008\%  \\ \hline
		Pirque       & 0,307\%  & 0,049\%  & 0,023\%  & 0,007\%  \\ \hline
		Puente Alto  & 0,328\%  & 0,056\%  & 0,023\%  & 0,008\%  \\ \hline
		Vitacura     & 0,355\%  & 0,057\%  & 0,027\%  & 0,000\%  \\ \hline
	\end{tabular}
\end{table}

%============================ Tabla 11 (abs6) transpuesta ==========================
\begin{table}[H]
	\centering
	\caption{Pérdida anual esperada (AAL) para fuente Intraplaca, usando modelo de fragilidad Cabrera et al. (2024).}
	\label{tab: AAL Intraplaca Cabrera 2024}
	\begin{tabular}{|c|c|c|c|c|}
		\hline
		Comuna       & \hdrunit{Adobe}{CLP $\cdot 10^{6}$} & \hdrunit{RC13}{CLP $\cdot 10^{6}$} & \hdrunit{RM}{CLP $\cdot 10^{6}$} & \hdrunit{Madera}{CLP $\cdot 10^{6}$} \\ \hline
		La Florida   & 0,000  & 0,000   & 0,119   & 0,044   \\ \hline
		La Reina     & 0,717  & 4,609   & 11,555  & 0,689   \\ \hline
		Las Condes   & 1,534  & 75,812  & 64,338  & 0,584   \\ \hline
		Peñalolén    & 0,000  & 1,491   & 4,472   & 0,068   \\ \hline
		Pirque       & 0,569  & 0,124   & 0,364   & 0,056   \\ \hline
		Puente Alto  & 0,766  & 1,901   & 2,937   & 0,173   \\ \hline
		Vitacura     & 1,183  & 0,858   & 11,830  & 0,000   \\ \hline
	\end{tabular}
\end{table}

%============================ Tabla 12 (perc6) transpuesta ==========================
\begin{table}[H]
	\centering
	\caption{Pérdida anual esperada relativa al total por comuna y por materialidad (AALR) para fuente Intraplaca, usando modelo de fragilidad Cabrera et al. (2024).}
	\label{tab: AALR Intraplaca Cabrera 2024}
	\begin{tabular}{|c|c|c|c|c|}
		\hline
		Comuna       & Adobe    & RC13     & RM       & Madera   \\ \hline
		La Florida   & 0,000\%  & 0,000\%  & 0,013\%  & 0,004\%  \\ \hline
		La Reina     & 0,286\%  & 0,040\%  & 0,018\%  & 0,006\%  \\ \hline
		Las Condes   & 0,247\%  & 0,038\%  & 0,018\%  & 0,006\%  \\ \hline
		Peñalolén    & 0,000\%  & 0,040\%  & 0,018\%  & 0,006\%  \\ \hline
		Pirque       & 0,247\%  & 0,037\%  & 0,018\%  & 0,006\%  \\ \hline
		Puente Alto  & 0,261\%  & 0,039\%  & 0,018\%  & 0,007\%  \\ \hline
		Vitacura     & 0,243\%  & 0,038\%  & 0,019\%  & 0,000\%  \\ \hline
	\end{tabular}
\end{table}

\subsubsection{Resultados Favier et al. (2017)}

%============================ Tabla 7 (abs4) transpuesta ==========================
\begin{table}[H]
	\centering
	\caption{Pérdida anual esperada (AAL) para fuente FSR, usando modelo de fragilidad Favier et al. (2017).}
	\label{tab: AAL FSR Favier 2017}
	\begin{tabular}{|c|c|c|c|}
		\hline
		Comuna       & \hdrunit{RC1024}{CLP $\cdot 10^{6}$} & \hdrunit{RC13}{CLP $\cdot 10^{6}$} & \hdrunit{RC49}{CLP $\cdot 10^{6}$} \\ \hline
		La Reina     & 0,000  & 1,057   & 0,000   \\ \hline
		Las Condes   & 0,663  & 17,500  & 10,000  \\ \hline
		Peñalolén    & 0,000  & 0,331   & 1,748   \\ \hline
		Pirque       & 0,000  & 0,029   & 0,000   \\ \hline
		Puente Alto  & 0,000  & 0,433   & 0,000   \\ \hline
		Vitacura     & 12,628 & 0,190   & 0,766   \\ \hline
	\end{tabular}
\end{table}

%============================ Tabla 8 (perc4) transpuesta ==========================
\begin{table}[H]
	\centering
	\caption{Pérdida anual esperada relativa al total por comuna y por materialidad (AALR) para fuente FSR, usando modelo de fragilidad Favier et al. (2017).}
	\label{tab: AALR FSR Favier 2017}
	\begin{tabular}{|c|c|c|c|}
		\hline
		Comuna       & RC1024    & RC13     & RC49     \\ \hline
		La Reina     & 0,000\%  & 0,009\%  & 0,000\%  \\ \hline
		Las Condes   & 0,012\%  & 0,009\%  & 0,011\%  \\ \hline
		Peñalolén    & 0,000\%  & 0,009\%  & 0,012\%  \\ \hline
		Pirque       & 0,000\%  & 0,008\%  & 0,000\%  \\ \hline
		Puente Alto  & 0,000\%  & 0,009\%  & 0,000\%  \\ \hline
		Vitacura     & 0,012\%  & 0,009\%  & 0,011\%  \\ \hline
	\end{tabular}
\end{table}



%============================ Tabla 9 (abs5) transpuesta ==========================
\begin{table}[H]
	\centering
	\caption{Pérdida anual esperada (AAL) para fuente Interplaca, usando modelo de fragilidad Favier et al. (2017).}
	\label{tab: AAL Interplaca Favier 2017}
	\begin{tabular}{|c|c|c|c|}
		\hline
		Comuna       & \hdrunit{RC1024}{CLP $\cdot 10^{6}$} & \hdrunit{RC13}{CLP $\cdot 10^{6}$} & \hdrunit{RC49}{CLP $\cdot 10^{6}$} \\ \hline
		La Reina     & 0,000  & 23,834  & 0,000   \\ \hline
		Las Condes   & 28,340 & 399,264 & 371,631 \\ \hline
		Peñalolén    & 0,000  & 7,760   & 68,978  \\ \hline
		Pirque       & 0,000  & 0,594   & 0,000   \\ \hline
		Puente Alto  & 0,000  & 10,153  & 0,000   \\ \hline
		Vitacura     & 537,731 & 4,668  & 31,241  \\ \hline
	\end{tabular}
\end{table}

%============================ Tabla 10 (perc5) transpuesta ==========================
\begin{table}[H]
	\centering
	\caption{Pérdida anual esperada relativa al total por comuna y por materialidad (AALR) para fuente Interplaca, usando modelo de fragilidad Favier et al. (2017).}
	\label{tab: AALR Interplaca Favier 2017}
	\begin{tabular}{|c|c|c|c|}
		\hline
		Comuna       & RC1024    & RC13     & RC49     \\ \hline
		La Reina     & 0,000\%  & 0,209\%  & 0,000\%  \\ \hline
		Las Condes   & 0,495\%  & 0,204\%  & 0,437\%  \\ \hline
		Peñalolén    & 0,000\%  & 0,207\%  & 0,462\%  \\ \hline
		Pirque       & 0,000\%  & 0,176\%  & 0,000\%  \\ \hline
		Puente Alto  & 0,000\%  & 0,213\%  & 0,000\%  \\ \hline
		Vitacura     & 0,494\%  & 0,211\%  & 0,464\%  \\ \hline
	\end{tabular}
\end{table}

%============================ Tabla 11 (abs6) transpuesta ==========================
\begin{table}[H]
	\centering
	\caption{Pérdida anual esperada (AAL) para fuente Intraplaca, usando modelo de fragilidad Favier et al. (2017).}
	\label{tab: AAL Intraplaca Favier 2017}
	\begin{tabular}{|c|c|c|c|}
		\hline
		Comuna       & \hdrunit{RC1024}{CLP $\cdot 10^{6}$} & \hdrunit{RC13}{CLP $\cdot 10^{6}$} & \hdrunit{RC49}{CLP $\cdot 10^{6}$} \\ \hline
		La Reina     & 0,000   & 16,034  & 0,000   \\ \hline
		Las Condes   & 17,900  & 262,990 & 255,304 \\ \hline
		Peñalolén    & 0,000   & 5,291   & 47,453  \\ \hline
		Pirque       & 0,000   & 0,488   & 0,000   \\ \hline
		Puente Alto  & 0,000   & 6,572   & 0,000   \\ \hline
		Vitacura     & 352,292 & 2,843   & 20,403  \\ \hline
	\end{tabular}
\end{table}

%============================ Tabla 12 (perc6) transpuesta ==========================
\begin{table}[H]
	\centering
	\caption{Pérdida anual esperada relativa al total por comuna y por materialidad (AALR) para fuente Intraplaca, usando modelo de fragilidad Favier et al. (2017).}
	\label{tab: AALR Intraplaca Favier 2017}
	\begin{tabular}{|c|c|c|c|}
		\hline
		Comuna       & RC1024    & RC13     & RC49     \\ \hline
		La Reina     & 0,000\%  & 0,137\%  & 0,000\%  \\ \hline
		Las Condes   & 0,313\%  & 0,134\%  & 0,298\%  \\ \hline
		Peñalolén    & 0,000\%  & 0,141\%  & 0,318\%  \\ \hline
		Pirque       & 0,000\%  & 0,140\%  & 0,000\%  \\ \hline
		Puente Alto  & 0,000\%  & 0,137\%  & 0,000\%  \\ \hline
		Vitacura     & 0,324\%  & 0,130\%  & 0,303\%  \\ \hline
	\end{tabular}
\end{table}


\subsubsection{Resultados HAZUS (2024)}


\begin{table}[H]
	\centering
	\caption{Pérdida anual esperada (AAL) para fuente FSR, usando modelo de fragilidad HAZUS (2024).}
	\label{tab: AAL FSR HAZUS 2024}
	\begin{tabular}{|c|c|c|c|c|c|}
		\hline
		Comuna       & \hdrunit{RC13}{CLP $\cdot 10^{6}$} & \hdrunit{RC47}{CLP $\cdot 10^{6}$} & \hdrunit{RC8}{CLP $\cdot 10^{6}$}  & \hdrunit{RM}{CLP $\cdot 10^{6}$}   & \hdrunit{Madera}{CLP $\cdot 10^{6}$} \\ \hline
		La Florida   & 0,000  & 0,000   & 0,000    & 0,045    & 0,038    \\ \hline
		La Reina     & 0,757  & 0,000   & 0,000    & 4,286    & 0,587    \\ \hline
		Las Condes   & 12,563 & 5,627   & 0,408    & 23,457   & 0,492    \\ \hline
		Peñalolén    & 0,237  & 1,024   & 0,000    & 1,649    & 0,059    \\ \hline
		Pirque       & 0,021  & 0,000   & 0,000    & 0,127    & 0,042    \\ \hline
		Puente Alto  & 0,311  & 0,000   & 0,000    & 1,007    & 0,129    \\ \hline
		Vitacura     & 0,137  & 0,443   & 7,761    & 4,054    & 0,000    \\ \hline
	\end{tabular}
\end{table}

%\newcommand{\hdrunit}[2]{\shortstack{#1\\(#2)}}

\begin{table}[H]
	\centering
	\caption{Pérdida anual esperada relativa al total por comuna y por materialidad (AALR) para fuente FSR, usando modelo de fragilidad HAZUS (2024).}
	\label{tab: AALR FSR HAZUS 2024}
	\begin{tabular}{|c|c|c|c|c|c|}
		\hline
		Comuna       & RC13    & RC47     & RC8      & RM      & Madera  \\ \hline
		La Florida   & 0,000\% & 0,000\%  & 0,000\%  & 0,005\% & 0,004\% \\ \hline
		La Reina     & 0,006\% & 0,000\%  & 0,000\%  & 0,007\% & 0,005\% \\ \hline
		Las Condes   & 0,006\% & 0,006\%  & 0,007\%  & 0,006\% & 0,005\% \\ \hline
		Peñalolén    & 0,006\% & 0,007\%  & 0,000\%  & 0,007\% & 0,005\% \\ \hline
		Pirque       & 0,006\% & 0,000\%  & 0,000\%  & 0,006\% & 0,005\% \\ \hline
		Puente Alto  & 0,006\% & 0,000\%  & 0,000\%  & 0,006\% & 0,005\% \\ \hline
		Vitacura     & 0,006\% & 0,007\%  & 0,007\%  & 0,006\% & 0,000\% \\ \hline
	\end{tabular}
\end{table}


%% =======================================================================
%% Conjunto 5
%% =======================================================================

\begin{table}[H]
	\centering
	\caption{Pérdida anual esperada (AAL) para fuente Interplaca, usando modelo de fragilidad HAZUS (2024).}
	\label{tab: AAL Interplaca HAZUS 2024}
	\begin{tabular}{|c|c|c|c|c|c|}
		\hline
		Comuna       & \hdrunit{RC13}{CLP $\cdot 10^{6}$} & \hdrunit{RC47}{CLP $\cdot 10^{6}$} & \hdrunit{RC8}{CLP $\cdot 10^{6}$}  & \hdrunit{RM}{CLP $\cdot 10^{6}$}   & \hdrunit{Madera}{CLP $\cdot 10^{6}$} \\ \hline
		La Florida   & 0,000  & 0,000    & 0,000    & 1,513    & 1,188    \\ \hline
		La Reina     & 28,933 & 0,000    & 0,000    & 165,332  & 20,408   \\ \hline
		Las Condes   & 469,620 & 267,668  & 24,906   & 916,890  & 17,733   \\ \hline
		Peñalolén    & 9,116  & 49,528   & 0,000    & 63,305   & 2,102    \\ \hline
		Pirque       & 0,711  & 0,000    & 0,000    & 4,326    & 1,301    \\ \hline
		Puente Alto  & 11,554 & 0,000    & 0,000    & 37,135   & 4,136    \\ \hline
		Vitacura     & 5,515  & 22,335   & 471,551  & 161,924  & 0,000    \\ \hline
	\end{tabular}
\end{table}

\begin{table}[H]
	\centering
	\caption{Pérdida anual esperada relativa al total por comuna y por materialidad (AALR) para fuente Interplaca, usando modelo de fragilidad HAZUS (2024).}
	\label{tab: AALR Interplaca HAZUS 2024}
	\begin{tabular}{|c|c|c|c|c|c|}
		\hline
		Comuna       & RC13    & RC47     & RC8      & RM      & Madera  \\ \hline
		La Florida   & 0,000\% & 0,000\%  & 0,000\%  & 0,166\% & 0,120\% \\ \hline
		La Reina     & 0,247\% & 0,000\%  & 0,000\%  & 0,256\% & 0,185\% \\ \hline
		Las Condes   & 0,239\% & 0,312\%  & 0,435\%  & 0,255\% & 0,187\% \\ \hline
		Peñalolén    & 0,243\% & 0,332\%  & 0,000\%  & 0,253\% & 0,184\% \\ \hline
		Pirque       & 0,208\% & 0,000\%  & 0,000\%  & 0,216\% & 0,149\% \\ \hline
		Puente Alto  & 0,242\% & 0,000\%  & 0,000\%  & 0,225\% & 0,167\% \\ \hline
		Vitacura     & 0,248\% & 0,332\%  & 0,433\%  & 0,254\% & 0,000\% \\ \hline
	\end{tabular}
\end{table}


%% =======================================================================
%% Conjunto 6
%% =======================================================================

\begin{table}[H]
	\centering
	\caption{Pérdida anual esperada (AAL) para fuente Intraplaca, usando modelo de fragilidad HAZUS (2024).}
	\label{tab: AAL Intraplaca HAZUS 2024}
	\begin{tabular}{|c|c|c|c|c|c|}
		\hline
		Comuna       & \hdrunit{RC13}{CLP $\cdot 10^{6}$} & \hdrunit{RC47}{CLP $\cdot 10^{6}$} & \hdrunit{RC8}{CLP $\cdot 10^{6}$}  & \hdrunit{RM}{CLP $\cdot 10^{6}$}   & \hdrunit{Madera}{CLP $\cdot 10^{6}$} \\ \hline
		La Florida   & 0,000   & 0,000   & 0,000    & 1,175    & 0,918    \\ \hline
		La Reina     & 21,947  & 0,000   & 0,000    & 124,632  & 15,405   \\ \hline
		Las Condes   & 354,634 & 206,058 & 18,390   & 671,593  & 12,874   \\ \hline
		Peñalolén    & 6,922   & 40,826  & 0,000    & 48,354   & 1,603    \\ \hline
		Pirque       & 0,537   & 0,000   & 0,000    & 3,385    & 1,023    \\ \hline
		Puente Alto  & 8,985   & 0,000   & 0,000    & 28,505   & 3,215    \\ \hline
		Vitacura     & 3,746   & 15,439  & 357,794  & 111,373  & 0,000    \\ \hline
	\end{tabular}
\end{table}

\begin{table}[H]
	\centering
	\caption{Pérdida anual esperada relativa al total por comuna y por materialidad (AALR) para fuente Interplaca, usando modelo de fragilidad HAZUS (2024).}
	\label{tab: AALR Intraplaca HAZUS 2024}
	\begin{tabular}{|c|c|c|c|c|c|}
		\hline
		Comuna       & RC13    & RC47     & RC8      & RM      & Madera  \\ \hline
		La Florida   & 0,000\% & 0,000\%  & 0,000\%  & 0,129\% & 0,092\% \\ \hline
		La Reina     & 0,187\% & 0,000\%  & 0,000\%  & 0,194\% & 0,139\% \\ \hline
		Las Condes   & 0,182\% & 0,239\%  & 0,321\%  & 0,185\% & 0,135\% \\ \hline
		Peñalolén    & 0,185\% & 0,273\%  & 0,000\%  & 0,193\% & 0,140\% \\ \hline
		Pirque       & 0,158\% & 0,000\%  & 0,000\%  & 0,169\% & 0,117\% \\ \hline
		Puente Alto  & 0,190\% & 0,000\%  & 0,000\%  & 0,173\% & 0,130\% \\ \hline
		Vitacura     & 0,170\% & 0,229\%  & 0,328\%  & 0,175\% & 0,000\% \\ \hline
	\end{tabular}
\end{table}

\end{appendixs}